% ============================================================
% Unconditional Negativity of the Second-Moment Bias
% and Off-Line Robustness
%
% Author:   Ulrich Tehrani
% Date:     June 2026
% Version:  v0.23
% Zenodo:   10.5281/zenodo.20792123
% License:  CC BY 4.0
% ============================================================

\documentclass[11pt]{article}
\usepackage{cmap}
\usepackage[T1]{fontenc}
\usepackage[utf8]{inputenc}
\usepackage[a4paper,margin=1in]{geometry}
\usepackage{amsmath,amssymb,amsthm,mathtools}
\usepackage{xcolor}
\usepackage{hyperref}
\usepackage{enumitem}
\usepackage{setspace}
\usepackage{booktabs}
\usepackage{array}
\usepackage{graphicx}

\hypersetup{colorlinks=true,
  linkcolor=blue!50!black,
  urlcolor=blue!50!black,
  citecolor=blue!50!black,
  pdftitle={Unconditional Negativity of the Second-Moment Bias and Off-Line Robustness},
  pdfauthor={Ulrich Tehrani}}
\setlength{\parskip}{0.5em}
\setlength{\parindent}{0pt}
\onehalfspacing
\emergencystretch=1em

% --------------------------------------------------------
% Theorem environments (numbered within section)
% --------------------------------------------------------
\newtheorem{theorem}{Theorem}[section]
\newtheorem{proposition}[theorem]{Proposition}
\newtheorem{lemma}[theorem]{Lemma}
\newtheorem{corollary}[theorem]{Corollary}
\newtheorem{definition}[theorem]{Definition}
\newtheorem{assumption}[theorem]{Assumption}
\newtheorem{observation}[theorem]{Observation}
\newtheorem{conjecture}[theorem]{Conjecture}
\newtheorem*{remark}{Remark}
\newtheorem{openproblem}[theorem]{Open Problem}

% --------------------------------------------------------
% Series-wide macros
% --------------------------------------------------------
\newcommand{\Hstr}{H_{\mathrm{str}}}
\newcommand{\Hnull}{H_{\mathrm{null}}}
\newcommand{\Tt}{\widetilde{T}}
\newcommand{\DSEL}{D_{\mathrm{SEL}}}

% --------------------------------------------------------
% Paper-8-specific macros
% --------------------------------------------------------
\newcommand{\BinfGW}{B_{\mathrm{GW}}^{\infty}}   % full Guinand--Weil second-moment curvature sum
\newcommand{\Binfline}{B_{\mathrm{line}}^{\infty}} % on-line (ordinate) proxy curvature sum
\newcommand{\dmin}{\delta_{\min}}           % binding near-zone gap
\newcommand{\ZtwoGW}{Z_{p,\mathrm{GW}}^{(2)}}   % full GW second-moment zero side
\newcommand{\Ztwoline}{Z_{p,\mathrm{line}}^{(2)}} % on-line ordinate proxy
\newcommand{\Pp}{P_p^{(2)}}                 % prime-side second-moment term
\newcommand{\Gp}{\Gamma_p^{(2)}}            % archimedean second-moment term
\newcommand{\Eoff}{E_p^{\mathrm{off},(2)}}  % off-line second-moment correction
\newcommand{\psiR}{\psi_{R}}                % Re psi on the quarter line

% --------------------------------------------------------
% Title
% --------------------------------------------------------
\title{\vspace{-1.5em}\bfseries Unconditional Negativity of the
Second-Moment Bias and Off-Line Robustness}
\author{Ulrich Tehrani}
\date{Research Note \textperiodcentered\ August~2026
      \textperiodcentered\ v0.23}

% ============================================================
\begin{document}
\maketitle
% ============================================================

\begin{abstract}
Building on the curvature identity $O''(\tfrac12)=-2B$ of~\cite{Paper6}
and the spectral trace formula $\mathrm{Tr}(\Tt(\sigma))=\DSEL-O(\sigma)$
of~\cite{Paper5}, this paper addresses, \emph{unconditionally}, two of the
open problems left by~\cite{Paper7}. Paper~7 left open a transfer from the
integrated bias to the direct second-moment curvature sum; the present paper
does \emph{not} prove that implication. Instead it \emph{bypasses} the
bridge, proving the direct second-moment negativity independently from the
differentiated Guinand--Weil formula, and it controls hypothetical
off-critical zeros.
We study the second Gaussian moment of the zero sum. Two objects must be
distinguished, the \emph{full Guinand--Weil} second-moment sum (over all
non-trivial zeros) and its \emph{on-line ordinate proxy},
\[
\ZtwoGW(\varepsilon)=\tfrac12\sum_\rho h^{(2)}_{p,\varepsilon}((\rho-\tfrac12)/i),
\qquad
\Ztwoline(\varepsilon)=\sum_{\gamma>0}e^{-\varepsilon^2\gamma^2}\gamma^2\cos(\gamma\log p);
\]
they coincide when all zeros lie on the critical line and otherwise differ by
an off-line correction $\Eoff(\varepsilon)=\ZtwoGW-\Ztwoline$.
With $\BinfGW(\kappa,\varepsilon)=\sum_{p\le\kappa}(\log p)^2\,\ZtwoGW(\varepsilon)$
and $\Binfline$ defined from the proxy, differentiating the Gaussian
explicit-formula identity in the smoothing parameter yields a second-moment
identity in which the diagonal prime self-interaction dominates; from it we
prove $\BinfGW(53,\varepsilon)<0$ for every $0<\varepsilon\le 0.05$
(closed-form for $\varepsilon\le 0.004$, certified by interval arithmetic up
to the reference width $\varepsilon=0.05$), and a uniform statement
$\BinfGW(\kappa,\varepsilon)<0$ for every $\kappa\ge 202$ and
$0<\varepsilon\le\varepsilon_0(\kappa)$ with the explicit threshold
$\varepsilon_0(\kappa)=1/(4e\kappa\sqrt{\log\kappa})$.
These sign statements are about the full Guinand--Weil object and assume
nothing about zero locations.
To pass from $\BinfGW$ to the operator curvature $O''(\tfrac12)=-2\Binfline$,
which is built from the on-line ordinates, we bound the off-line correction:
a $\beta$-uniform magnitude estimate combined with the Platt--Trudgian
verified height $H_0=3\cdot10^{12}$~\cite{PlattTrudgian21} shows that
$\sum_{p\le\kappa}(\log p)^2|\Eoff|$ stays below the certified negativity
margin $M$ for every smoothing width $\varepsilon\ge\varepsilon_{\mathrm{off}}(\kappa,M,H_0)$;
within a certified full-object negativity range this transfers the sign, so
for $\varepsilon_{\mathrm{off}}(\kappa,M,H_0)\le\varepsilon\le 0.05$ (in particular
$\varepsilon=0.05$) one has $\Binfline(53,\varepsilon)<0$ and
$O''(\tfrac12)>0$.
No hypothetical input is used: no Riemann Hypothesis, no Hilbert--P\'olya
postulate, no random-matrix conjecture; the off-line bound rests on an
unconditional finite-height verification, not on RH.

\textbf{What is proved.}
The second-moment explicit-formula identity (Theorem~\ref{thm:id});
the archimedean bound $\Gp=O(1)$ (Theorem~\ref{thm:arch});
prime-side eigenterm extraction and cross-prime control
(Theorem~\ref{thm:prime});
$\BinfGW(53,\varepsilon)<0$ for $0<\varepsilon\le 0.05$
(Theorem~\ref{thm:ref});
the uniform negativity threshold $\varepsilon_0(\kappa)$ for $\kappa\ge 202$
(Theorem~\ref{thm:uniform}); off-line robustness in magnitude form
(Theorem~\ref{thm:offline}); and the operator curvature
$O''(\tfrac12)=-2\Binfline>0$ for
$\varepsilon_{\mathrm{off}}\le\varepsilon\le 0.05$,
in particular at $\varepsilon=0.05$ (Corollary~\ref{cor:curvline}).

\textbf{What is numerical.}
At $(\kappa,\varepsilon,N)=(53,0.05,100)$:
$B(53,0.05,100)=-19\,342.5$; $O''(\tfrac12)=+38\,685.1$;
$S_3(53)=87.4247$; $\dmin(53)=0.031749$.
The archimedean bound $I_4(\tfrac14)\le 3561.1$ is \emph{proved}
(Theorem~\ref{thm:arch}) and \emph{certified} in Arb ball arithmetic
(\texttt{certify\_I4}).

\textbf{What remains open.}
The Cancellation Hypothesis~(A$^{**}$) of~\cite{Paper7};
the Weighted Bias Bridge implication
$B_{\mathrm{int}}^\infty<0\Rightarrow\BinfGW<0$ (bypassed here, not proved);
exact stationarity $O'(\tfrac12)=0$;
the Weil-transfer operator and Weil positivity;
and the asymptotics of the negativity window.

\textbf{Scope.}
The negativity $\BinfGW<0$ of the full Guinand--Weil object, and
$\Binfline<0$ of the on-line proxy for
$\varepsilon_{\mathrm{off}}\le\varepsilon\le 0.05$,
are statements about the second-moment bias and the curvature
$O''(\tfrac12)=-2\Binfline>0$.
It is \emph{not} a criterion for the Riemann Hypothesis; this paper does
not claim and does not suggest RH (see the scope remark in
\S\,\ref{sec:open}).

\textbf{Structure.}
Sections~\ref{sec:id}--\ref{sec:prime} establish the second-moment
identity and the three side estimates;
Sections~\ref{sec:ref}--\ref{sec:uniform} prove negativity at the
reference parameter and uniformly in the cutoff;
Section~\ref{sec:offline} proves off-line robustness.
All results are unconditional (Theorem~\ref{thm:offline} is a magnitude
statement); no use of RH, GUE, or the Hilbert--P\'olya postulate is made.
\end{abstract}

\vspace{2em}

% --------------------------------------------------------
% Notation and Conventions
% --------------------------------------------------------
\section*{Notation and Conventions}
\label{sec:notation}

\emph{Critical line.}
Throughout, the critical line means $\Re(s)=\tfrac12$.
The trace parameter $\sigma\in\mathbb{R}$ is evaluated at $\sigma=\tfrac12$
unless stated otherwise.

\medskip
\emph{Global parameters.}
The reference values are $\kappa=53$ (prime cutoff, yielding $16$ active
primes $p\le\kappa$), $\varepsilon=0.05$ (Gaussian regularisation
parameter), and $N=100$ (number of Riemann ordinates $\gamma_k$ employed).
These values are fixed throughout unless explicitly varied.
We write $P(\kappa):=\sum_{p\le\kappa}\log p=\vartheta(\kappa)$ for the
Chebyshev theta sum and
$S_3(\kappa):=\sum_{p\le\kappa}(\log p)^3/\sqrt{p}$.
At the reference parameters $P(53)\approx 44.93$ and $S_3(53)=87.4247$.

\medskip
\emph{Main objects.}
The coupling operator $\Phi(\sigma):\Hstr\to\Hnull$ is defined by
\[
\Phi(\sigma)_{k,p}=e^{-\varepsilon^2\gamma_k^2/2}\sin(\sigma\gamma_k\log p),
\]
and the Tehrani operator is $\Tt(\sigma)=\Phi(\sigma)\circ\Phi(\sigma)^*$
(\cite{Paper5}).
The trace oscillatory sum is
\[
O(\sigma)=\tfrac12\sum_{k,p}e^{-\varepsilon^2\gamma_k^2}\cos(2\sigma\gamma_k\log p),
\]
with curvature identity $O''(\tfrac12)=-2B$ (\cite{Paper6}), where
\[
B=\sum_{k,p}e^{-\varepsilon^2\gamma_k^2}(\gamma_k\log p)^2\cos(\gamma_k\log p)
\]
is the finite on-line (ordinate) curvature sum.
The present paper works with two infinite-$N$ second-moment objects that
must be kept distinct. The \emph{full Guinand--Weil} second-moment zero side
and its \emph{on-line ordinate proxy} are
\[
\ZtwoGW(\varepsilon)
 := \tfrac12\sum_\rho h^{(2)}_{p,\varepsilon}\!\Bigl(\tfrac{\rho-\tfrac12}{i}\Bigr),
\qquad
\Ztwoline(\varepsilon)
 := \sum_{\substack{\rho=\tfrac12+i\gamma\\ \gamma>0}}
    e^{-\varepsilon^2\gamma^2}\,\gamma^2\cos(\gamma\log p),
\]
the sum in $\ZtwoGW$ ranging over all non-trivial zeros $\rho$ with
multiplicity; their difference is the off-line correction
$\Eoff(\varepsilon):=\ZtwoGW(\varepsilon)-\Ztwoline(\varepsilon)$, which
vanishes when all zeros lie on the critical line. The corresponding
curvature sums are
\[
\BinfGW(\kappa,\varepsilon)
 := \sum_{p\le\kappa}(\log p)^2\,\ZtwoGW(\varepsilon),
\qquad
\Binfline(\kappa,\varepsilon)
 := \sum_{p\le\kappa}(\log p)^2\,\Ztwoline(\varepsilon),
\]
so that $\Binfline$ is the infinite-$N$ form of $B$, and
$O''(\tfrac12)=-2\Binfline$.
At the reference parameters the finite-$N$ curvature sum is
$B(53,0.05,100)=-19\,342.5$, which agrees with $\Binfline(53,0.05)$ up to the
truncation control quoted from~\cite{Paper6}; we use it as the finite-$N$
numerical proxy for $\Binfline(53,0.05)$.
The Gaussian test functions are
\[
h^{(0)}_{p,s}(u)=e^{-s u^2}\cos(u\log p),
\qquad
h^{(2)}_{p,\varepsilon}(u)=u^2 e^{-\varepsilon^2 u^2}\cos(u\log p)
 = -\,\partial_s h^{(0)}_{p,s}(u)\big|_{s=\varepsilon^2},
\]
both even and entire.
We write $\Lambda$ for the von Mangoldt function and
$\psiR(t):=\Re\,\psi(\tfrac14+\tfrac{it}{2})$ for the real part of the
digamma function on the relevant vertical line; the archimedean
second-moment term is
$\Gp(\varepsilon)=\frac{1}{2\pi}\int_0^\infty t^2 e^{-\varepsilon^2 t^2}
\cos(t\log p)\,\psiR(t)\,dt$.

\medskip
\emph{Analytical foundations.}
The proofs use the Guinand--Weil explicit formula for Gaussian test
functions in the form of~\cite[Thm.~5.12]{IwaniecKowalski04} (the
historical sources are~\cite{Guinand,Weil}), the Riemann--von Mangoldt
zero-counting asymptotics
$N(T)\sim\frac{T}{2\pi}\log\frac{T}{2\pi}$ in the explicit form
of~\cite{HasanalizadeShenWong21}, and the Chebyshev-type bound
$\psi(x)\le 1.04\,x$ of~\cite{RosserSchoenfeld62}.
Section~\ref{sec:offline} uses, as a finite verified arithmetic fact, the
unconditional interval-arithmetic computation of
Platt and Trudgian~\cite{PlattTrudgian21}: every non-trivial zero
$\beta+i\gamma$ with $0<\gamma\le H_0=3\cdot10^{12}$ has $\beta=\tfrac12$.
No zero-location, pair-correlation, or spectral hypothesis is used.

\medskip
\emph{Numerical computations.}
The illustrative tables and calibration values of this paper (anchors such as
$B(53,0.05,100)=-19\,342.5$, the residual table, the $\kappa=202$ budgets) were
produced with \texttt{mpmath} at $50$~decimal digits and serve as candidate
values, not as certificates; Riemann zeros $\gamma_k$ were evaluated via
\texttt{mpmath.zetazero}.
The word \emph{certified} is reserved for the threshold statements of
\S\,\ref{sec:ref}--\ref{sec:uniform}, which are established by a separate
rigorous \emph{interval} computation in Arb ball arithmetic (midpoint--radius
with directed outward rounding)~\cite{Johansson17}, as provided by the
FLINT~3 library (Arb layer) through its Python bindings, with proven Gaussian
tail bounds below $10^{-380}$. The certificate recomputes every prime-power
interval contribution, bounds the discarded tail, and prints lower bound,
right-hand side, and margin for each covering interval, aborting on any
overlap, gap, or non-positive margin.
Verification code (the consistency gate and the interval certificate) is
available at~\cite{Code}
(\url{https://github.com/utehrani/analysislab-nt}).

\medskip
\emph{Not to be confused with.}
The following symbols look similar but are distinct objects.
\begin{description}[leftmargin=2.6cm,style=nextline,itemsep=2pt]
\item[$\ZtwoGW$ vs.\ $\Ztwoline$] full Guinand--Weil zero side
  $\tfrac12\sum_\rho$ versus the on-line ordinate proxy $\sum_{\gamma>0}$;
  their difference is $\Eoff=\ZtwoGW-\Ztwoline$.
\item[$\BinfGW$ vs.\ $\Binfline$] full-object sign (this paper) versus the
  proxy curvature sum (the finite-$N$ ordinate computation gives $-19\,342.5$
  at $N=100$, consistent with the sign obtained after the off-line transfer).
\item[$\BinfGW,\Binfline$ vs.\ $B_{\mathrm{int}}$] the $\gamma_k^2$-weighted
  curvature sums versus the integrated bias: the finite reference value is
  $B_{\mathrm{int}}^+(0.05,100)=-42.21$ (\cite{Paper6,Paper7}), distinct from
  the infinite integrated-bias window $B_{\mathrm{int}}^\infty$ (\cite{Paper7}).
\item[pole $-\tfrac18$ vs.\ ratio $-\tfrac14$] the pole coefficient of
  $2e^{\varepsilon^2/4}\cosh(\tfrac{\log p}{2})$ is $-\tfrac18$; the ratio to
  the zeroth-moment pole term is $-\tfrac14$.
\item[$\varepsilon_0(\kappa)$ vs.\ $\varepsilon_{\mathrm{int}}(\kappa)$] the
  uniform threshold $1/(4e\kappa\sqrt{\log\kappa})$ of this paper versus the
  integrated-window scale of~\cite{Paper7}.
\item[$\varepsilon\le 0.004$ vs.\ $\varepsilon\le 0.05$] the closed-form
  threshold versus the certified threshold.
\end{description}
Both pole numbers are correct: the pole term of the second-moment identity
equals $-\tfrac18\cdot 2e^{\varepsilon^2/4}\cosh(\tfrac{\log p}{2})$, i.e.\
exactly minus one quarter of the zeroth-moment pole term.
The closed-form and certified thresholds in Theorem~\ref{thm:ref} are
stated separately and never conflated.

\medskip
\emph{Object summary.}
\begin{center}
\begin{tabular}{lll}
\toprule
Object & Definition & Role \\
\midrule
$\ZtwoGW$ & $\tfrac12\sum_\rho h^{(2)}_{p,\varepsilon}((\rho-\tfrac12)/i)$ & full GW zero side \\
$\Ztwoline$ & $\sum_{\gamma>0} e^{-\varepsilon^2\gamma^2}\gamma^2\cos(\gamma\log p)$ & on-line proxy \\
$\Eoff$ & $\ZtwoGW-\Ztwoline$ & off-line correction \\
$\BinfGW$ & $\sum_{p\le\kappa}(\log p)^2\,\ZtwoGW$ & proved negative \\
$\Binfline$ & $\sum_{p\le\kappa}(\log p)^2\,\Ztwoline$ & operator curvature \\
$B(\kappa,\varepsilon,N)$ & finite ordinate truncation & numerical proxy \\
\bottomrule
\end{tabular}
\end{center}

\vspace{2em}
\newpage
% ============================================================
\section{Introduction}
\label{sec:intro}
% ============================================================

\subsection{Background and Motivation}

This paper is the eighth in a series developing an operator-theoretic
framework for studying the distribution of the non-trivial zeros of the
Riemann zeta function, in which prime numbers mediate the coupling between
zeros through an explicit family of finite-dimensional operators.

\textbf{Paper~1}~\cite{Paper1} introduced a local curvature decomposition
of the explicit formula and identified the critical line as a phase
boundary through a structural divergence of the curvature at
$\sigma=\tfrac12$.

\textbf{Paper~2}~\cite{Paper2} connected this local curvature to the Weil
explicit functional through $\sigma$-dependent admissible Gaussian test
functions, and identified the Weil positivity framework
of~\cite{Lagarias} as the natural target.

\textbf{Paper~3}~\cite{Paper3} made the geometric structure explicit by
introducing two finite-dimensional Hilbert spaces, $\Hstr$ over primes and
$\Hnull$ over zeros, together with the coupling operator
$\Phi:\Hstr\to\Hnull$ built from prime--zero phasors, and the loop
operator $T=\Phi^*\circ\Phi$.

\textbf{Paper~4}~\cite{Paper4} introduced the dual Tehrani operator
$\Tt=\Phi\circ\Phi^*$ on $\Hnull$ and showed, through a
Hilbert--P\'olya diagnostic ($r_2\to 0.16$ as $\kappa\to\infty$
[numerical]), that $\Tt$ is \emph{not} a Hilbert--P\'olya operator,
distinguishing it from the constructions of Berry and
Keating~\cite{BerryKeating} and de Branges~\cite{deBranges}.

\textbf{Paper~5}~\cite{Paper5} extended the operator family to a
one-parameter family $\Tt(\sigma)=\Phi(\sigma)\circ\Phi(\sigma)^*$ and
proved the exact spectral trace identity
$\mathrm{Tr}(\Tt(\sigma))=\DSEL-O(\sigma)$, where $\DSEL$ is
$\sigma$-independent.
It introduced the Guinand--Weil smoothed zero sum as motivation for the
bias of prime--zero sums.

\textbf{Paper~6}~\cite{Paper6} proved the curvature--bias identity
$O''(\tfrac12)=-2B$ by direct differentiation of the trace formula, with
$B=\sum_{k,p}e^{-\varepsilon^2\gamma_k^2}(\gamma_k\log p)^2
\cos(\gamma_k\log p)$, and recorded the reference value $B=-19\,342.5$,
giving $O''(\tfrac12)=+38\,685.1>0$ [numerical].
The transfer from the integrated bias to the direct curvature sum was left
open; the present paper bypasses it (see below).

\textbf{Paper~7}~\cite{Paper7} proved, under two named hypotheses, the
pointwise asymptotic rate $r_p^\infty=\tfrac12$ and a conditional window
$B_{\mathrm{int}}^\infty(\kappa,\varepsilon)<0$, and formulated the
\emph{Weighted Bias Bridge} --- the transfer from the integrated bias
$B_{\mathrm{int}}^\infty$ to the direct second-moment sum $\BinfGW$ --- and
the question of off-critical control as its two central open problems.

\smallskip
\textbf{The present paper} addresses both of these open problems
\emph{unconditionally}, by bypassing the bridge rather than establishing it:
we do not prove the implication
$B_{\mathrm{int}}^\infty<0\Rightarrow\BinfGW<0$, but prove $\BinfGW<0$
directly.
We work directly with the second-moment object $\BinfGW(\kappa,\varepsilon)$,
defined through the zero side of the Guinand--Weil explicit formula and
hence free of any zero-location hypothesis.
The mechanism is elementary in outline: differentiating the Gaussian
explicit-formula identity once in the smoothing parameter $s=\varepsilon^2$
produces a second-moment identity whose dominant term, for small
$\varepsilon$, is the diagonal self-interaction of each prime with its own
term in the explicit formula.
This term is negative and of size $\varepsilon^{-3}$; every competing
contribution --- the pole term, the cross-prime interactions, and the
archimedean integral --- is of strictly lower order, with explicit
constants.
The negativity of $\BinfGW$ for small $\varepsilon$ is therefore an
unconditional, purely prime-side phenomenon.

The operator $\Tt(\sigma)$ of this series is a Gram-type operator with
explicit arithmetic matrix entries; it is not a Hilbert--P\'olya operator
(\cite{Paper4}), and does not arise from an adelic or non-commutative
geometric construction.
Connes and Consani~\cite{ConnesConsani23} and Connes, Consani, and
Moscovici~\cite{CCM25} construct spectral triples and zeta-cycles on
adelic spaces; our framework is finite-cutoff, explicit, and classical.
Here $\BinfGW$ is an analytic bias object --- the second derivative in
$\sigma$ of the spectral trace --- not an adelic structure.
The methodology of replacing a zero-location hypothesis by a certified
finite-height verification combined with explicit-formula estimates is
itself well established; the closest mean-square neighbour is the work of
Brent, Platt, and Trudgian~\cite{BrentPlattTrudgian20} on the error term
in the prime number theorem, and related Hilbert-space approaches to the
Weil distribution have been pursued by Burnol~\cite{Burnol} and
Suzuki~\cite{Suzuki25}.
The bridge to the Weil positivity criterion of~\cite{Lagarias,Bombieri}
remains the principal open problem of the series.

\subsection{Main Results}

The six main results of this paper are as follows.

The first, in \S\,\ref{sec:id} (Theorem~\ref{thm:id}), is the
\emph{second-moment identity}: an explicit decomposition of
$\ZtwoGW(\varepsilon)$ into a pole term with coefficient $-\tfrac18$
(equivalently, minus one quarter of the zeroth-moment pole term), a full
prime-power term over all $q^m$, and an archimedean integral, obtained by
differentiating the Gaussian explicit formula in $s=\varepsilon^2$.

The second, in \S\,\ref{sec:arch} (Theorem~\ref{thm:arch}), is the
\emph{archimedean bound} $\Gp(\varepsilon)=O(1)$, with an explicit constant
$I_4\le 3561.1$ obtained by four-fold integration by parts against
$\cos(t\log p)$ using $g'''(0)=0$; this term is subdominant by three powers
of $\varepsilon$.

The third, in \S\,\ref{sec:prime} (Theorem~\ref{thm:prime}), is the
\emph{prime-side control}: the diagonal $(q,m)=(p,1)$ term contributes
exactly $\log p/(\sqrt p\,8\sqrt\pi\,\varepsilon^3)$, and the cross-prime
remainder is controlled through the binding near-zone gap
$\dmin(53)=0.031749$ (attained by the pair $31\leftrightarrow 2^5$).

The fourth, in \S\,\ref{sec:ref} (Theorem~\ref{thm:ref}), is
\emph{negativity at the reference parameter}: $\BinfGW(53,\varepsilon)<0$ for
all $0<\varepsilon\le 0.05$.
The closed-form criterion proves this for $\varepsilon\le 0.004$; a
certification by interval arithmetic with the same proven constants extends
it to the reference width $\varepsilon=0.05$ (and, by chaining interval
certificates, to $\varepsilon\le 0.0819$).
These two stages are stated separately.
This certifies the sign of the full Guinand--Weil object; the operator
curvature $O''(\tfrac12)=-2\Binfline>0$, built from the on-line ordinates,
follows once the off-line correction is bounded
(Corollary~\ref{cor:curvline}, after Theorem~\ref{thm:offline}).

The fifth, in \S\,\ref{sec:uniform} (Theorem~\ref{thm:uniform}), is
\emph{uniform negativity}: for every cutoff $\kappa\ge 202$ and every
$0<\varepsilon\le\varepsilon_0(\kappa)$ with
$\varepsilon_0(\kappa)=1/(4e\kappa\sqrt{\log\kappa})$, one has
$\BinfGW(\kappa,\varepsilon)<0$.
This establishes the direct second-moment negativity that the Weighted Bias
Bridge would have yielded, but proves it directly rather than through the
bridge.

The sixth, in \S\,\ref{sec:offline} (Theorem~\ref{thm:offline}), is
\emph{off-line robustness} in magnitude form: a $\beta$-uniform bound on
the contribution of any hypothetical off-critical zeros, combined with the
Platt--Trudgian verified height $H_0=3\cdot10^{12}$ and an explicit discrete
zero-count inequality, shows that the off-line correction $\Eoff$ stays
below the certified negativity margin $M$ for
$\varepsilon\ge\varepsilon_{\mathrm{off}}$, without any RH assumption. Within
a certified negativity range this transfers the sign from $\BinfGW$ to
$\Binfline$, so for $\varepsilon_{\mathrm{off}}\le\varepsilon\le 0.05$ (in
particular $\varepsilon=0.05$) one has $O''(\tfrac12)>0$
(Corollary~\ref{cor:curvline}).

\paragraph{Relation to existing work.}
The Guinand--Weil explicit formula, its prime-power and archimedean
(digamma) components, the broader literature in which signs are read off
from zeros, and the general strategy of removing an RH assumption through a
certified finite-height zero verification are all classical; the
Gaussian/Hermite-weighted identity technology appears in nearby form in
recent work~\cite{BalanzarioCardenas23}. What is new here is the
conjunction: the unconditional pointwise sign $\BinfGW<0$ for this specific
Gaussian second-moment curvature functional, the constant-tracked
archimedean bound (four-fold integration by parts, $I_4\le 3561.1$)
tailored to the twice-differentiated identity, and the transfer of the
certified-zero methodology to off-line robustness of a curvature sign
statement. The identity of \S\,\ref{sec:id} is a specialized differentiated
descendant of the classical explicit formula, not a new formula; the
smoothed and weighted explicit-formula framework, the prime-power/digamma
decomposition, and the certified-zero method itself are cited here as prior
art, not claimed as contributions.

\subsection{Reader Contract}

This paper \emph{proves} (unconditionally):
the second-moment explicit-formula identity (Theorem~\ref{thm:id});
the archimedean bound $\Gp=O(1)$ (Theorem~\ref{thm:arch});
prime-side eigenterm extraction and cross-prime control
(Theorem~\ref{thm:prime});
$\BinfGW(53,\varepsilon)<0$ for all $0<\varepsilon\le 0.05$
(Theorem~\ref{thm:ref});
the uniform threshold $\BinfGW(\kappa,\varepsilon)<0$ for $\kappa\ge 202$,
$\varepsilon\le\varepsilon_0(\kappa)$ (Theorem~\ref{thm:uniform});
off-line robustness in magnitude form (Theorem~\ref{thm:offline});
and the operator curvature positivity $O''(\tfrac12)=-2\Binfline>0$ for
$\varepsilon_{\mathrm{off}}\le\varepsilon\le 0.05$, in particular at
$\varepsilon=0.05$ (Corollary~\ref{cor:curvline}).

At the reference parameters $(\kappa,\varepsilon,N)=(53,0.05,100)$ this paper
\emph{establishes numerically} $B(53,0.05,100)=-19\,342.5$,
$O''(\tfrac12)=+38\,685.1$, and $\dmin(53)=0.031749$.
It \emph{proves} the archimedean bound $I_4\le 3561.1$ and the closed-form
threshold $\varepsilon_0(53)=0.004$, and \emph{certifies} (interval
arithmetic) negativity for $\varepsilon\le 0.05$ with single-interval margins
$0.83$ on $(0,0.047]$ and $24.29$ on $[0.047,0.05]$.

This paper \emph{leaves open}:
the Cancellation Hypothesis~(A$^{**}$) of~\cite{Paper7};
the Weighted Bias Bridge implication
$B_{\mathrm{int}}^\infty<0\Rightarrow\BinfGW<0$ itself (the present paper
bypasses it, proving $\BinfGW<0$ directly, rather than establishing the
implication);
exact stationarity $O'(\tfrac12)=0$;
the Weil-transfer operator and Weil positivity;
and the asymptotics of the negativity window
$\inf_\kappa\varepsilon_{\mathrm{int}}(\kappa)$.

This paper does \emph{not} claim and does \emph{not} suggest the Riemann
Hypothesis.
The negativity $\BinfGW<0$ is a second-moment bias / curvature statement, not
an RH criterion: the canonical off-critical continuation of the curvature
test function yields no sign-definite $\delta^2$-defect, as recorded in the
scope remark of \S\,\ref{sec:open}.

% ============================================================
\section{Background: Trace Formula, Curvature Identity, and the
Guinand--Weil Framework}
\label{sec:bg}
% ============================================================

This short section records the inputs imported from the preceding papers;
all are cited, none re-proved here.

\subsection{Trace formula and curvature identity}

Paper~5~\cite{Paper5} proves the exact identity
$\mathrm{Tr}(\Tt(\sigma))=\DSEL-O(\sigma)$ with $\DSEL$ $\sigma$-independent
and
$O(\sigma)=\tfrac12\sum_{k,p}e^{-\varepsilon^2\gamma_k^2}
\cos(2\sigma\gamma_k\log p)$.
Differentiating twice in $\sigma$ and evaluating at $\sigma=\tfrac12$,
Paper~6~\cite{Paper6} obtains the curvature identity
\[
O''(\tfrac12) = -2B,
\qquad
B = \sum_{k,p}e^{-\varepsilon^2\gamma_k^2}(\gamma_k\log p)^2
   \cos(\gamma_k\log p).
\]
The infinite-$N$, on-line form of $B$ is
$\Binfline(\kappa,\varepsilon)=\sum_{p\le\kappa}(\log p)^2\Ztwoline(\varepsilon)$,
so that $O''(\tfrac12)=-2\Binfline$ and the sign of $\Binfline$ is the sign
of $-O''(\tfrac12)$.
The present paper proves negativity for the full Guinand--Weil object
$\BinfGW(\kappa,\varepsilon)=\sum_{p\le\kappa}(\log p)^2\ZtwoGW(\varepsilon)$
on the prime/pole/archimedean side (no zero-location hypothesis), and then
transfers the sign to $\Binfline$ by bounding $\Eoff$ in
\S\,\ref{sec:offline}.

\begin{proposition}[Infinite-$N$ curvature identity; proved]\label{prop:infN}
Write
\[
O_\infty(\sigma)=\tfrac12\sum_{k,p}e^{-\varepsilon^2\gamma_k^2}
\cos(2\sigma\gamma_k\log p)
\]
for the $N\to\infty$ on-line trace. For each fixed
$\varepsilon>0$ the series for $O_\infty$ and its first two $\sigma$-derivatives
converge locally uniformly in $\sigma$, so $O_\infty\in C^2$ and term-by-term
differentiation is valid; in particular
\[
O_\infty''(\tfrac12)=-2\Binfline(\kappa,\varepsilon).
\]
\end{proposition}

\begin{proof}
The $\nu$-th $\sigma$-derivative of the $(k,p)$ term is bounded by
$e^{-\varepsilon^2\gamma_k^2}(2\gamma_k\log p)^\nu$. By Riemann--von Mangoldt,
$N(T)=O(T\log T)$, so for fixed $p$ the ordinate sum
$\sum_k e^{-\varepsilon^2\gamma_k^2}\gamma_k^{\nu}$ is dominated by
\[
\int_0^\infty t^\nu e^{-\varepsilon^2 t^2}\,dN(t)
=O\!\Bigl(\int_0^\infty t^{\nu+1}\log t\,e^{-\varepsilon^2 t^2}\,dt\Bigr)<\infty,
\]
the Gaussian damping overpowering the $O(T\log T)$ growth for every fixed
$\varepsilon>0$ and $\nu\le 2$. The bound is independent of $\sigma$, so by the
Weierstrass $M$-test each of $O_\infty,O_\infty',O_\infty''$ converges
uniformly on compact $\sigma$-sets; differentiation term by term is therefore
legitimate, and evaluating the second derivative at $\sigma=\tfrac12$ gives
$-2\sum_{k,p}e^{-\varepsilon^2\gamma_k^2}(\gamma_k\log p)^2\cos(\gamma_k\log p)
=-2\Binfline$.
\end{proof}



\subsection{The Guinand--Weil framework and its second moment}

For an admissible even test function $h$ the Guinand--Weil explicit formula
\cite[Thm.~5.12]{IwaniecKowalski04} relates the zero side
$\tfrac12\sum_{\rho}h\bigl((\rho-\tfrac12)/i\bigr)$ (sum over all
non-trivial zeros $\rho$ with multiplicity) to a pole term, a prime-power
sum against the Fourier transform $\widehat h$, and an archimedean integral
against the digamma function.
When all zeros lie on the critical line the zero side equals
$\sum_{k\ge 1}h(\gamma_k)$; this is the form used in all numerical
computations, the zeros entering the computations ($\gamma_k\le 1002$) being
classically known to lie on the line.
\emph{No zero-location hypothesis enters any proof below}: all zero sums are
defined through the zero side of the explicit formula.

The \emph{second moment} arises from the second $\sigma$-derivative of the
trace, equivalently from differentiating the even Gaussian test function
$h^{(0)}_{p,s}(u)=e^{-s u^2}\cos(u\log p)$ once in $s$ at $s=\varepsilon^2$:
$h^{(2)}_{p,\varepsilon}=-\partial_s h^{(0)}_{p,s}|_{s=\varepsilon^2}$.
This is the object differentiated throughout \S\,\ref{sec:id}.

\begin{remark}
The full Guinand--Weil object and the ordinate proxy differ, when
off-critical zeros are present, by the off-line contribution
$\Eoff$ isolated in \S\,\ref{sec:offline}; under the all-zeros-on-line
identity of Paper~5 they coincide.
The present paper never assumes this identity: \S\S\,\ref{sec:id}--\ref{sec:uniform}
work on the prime/pole/archimedean side, and \S\,\ref{sec:offline} bounds
$\Eoff$ explicitly.
\end{remark}

% ============================================================
\section{The Second-Moment Identity}
\label{sec:id}
% ============================================================

We first record absolute convergence, then the explicit identity. The
differentiated Gaussian identity below belongs to the same family of
weighted explicit-formula identities as the Gaussian/Hermite-weighted zero
sums of~\cite{BalanzarioCardenas23}; that work is, however, conditional on
the Riemann Hypothesis and simplicity of zeros and aims at localization,
whereas the unconditional sign statement proved here is of a different
nature.

\begin{lemma}[Absolute convergence of the second moment; proved]
\label{lem:conv}
For every $\varepsilon>0$ and every prime $p$ the sum defining
$\ZtwoGW(\varepsilon)$ converges absolutely, and so does the full zero-side
sum $\tfrac12\sum_\rho h^{(2)}_{p,\varepsilon}\bigl((\rho-\tfrac12)/i\bigr)$.
\end{lemma}

\begin{proof}
Write $\rho=\beta+i\gamma$ with $0<\beta<1$ and
$z=(\rho-\tfrac12)/i=\gamma-i(\beta-\tfrac12)$.
Then $\Re(z^2)=\gamma^2-(\beta-\tfrac12)^2\ge\gamma^2-\tfrac14$ and
$|\cos(zL)|\le e^{|\Im z|L}\le\sqrt p$ with $L=\log p$, so
\[
\bigl|h^{(2)}_{p,\varepsilon}(z)\bigr|
 \le \bigl(\gamma^2+\tfrac14\bigr)\,e^{\varepsilon^2/4}\sqrt p\;
   e^{-\varepsilon^2\gamma^2}.
\]
By the Riemann--von Mangoldt formula in the explicit form
$N(T)=F(T)+O(\log T)$ with $F(T)=\tfrac{T}{2\pi}\log\tfrac{T}{2\pi}
-\tfrac{T}{2\pi}+\tfrac78$ and the quantitative remainder
$|N(T)-F(T)|\le E(T)=O(\log T)$ of~\cite{HasanalizadeShenWong21} (used
quantitatively in \S\,\ref{sec:offline}), the number of zeros with
$n\le|\gamma|<n+1$ is
\[
N(n+1)-N(n)\le \bigl(F(n+1)-F(n)\bigr)+E(n+1)+E(n)=O(\log(n+2)),
\]
since the differenced main term $F(n+1)-F(n)=\tfrac1{2\pi}\log\tfrac{n}{2\pi}
+O(1)$; hence the sum is dominated by
$\sum_{n\ge 0}(n+1)^2\log(n+2)e^{-\varepsilon^2 n^2}<\infty$.
The ordinate sum is the special case $\beta=\tfrac12$.
\end{proof}

\begin{theorem}[Second-moment explicit-formula identity; proved]
\label{thm:id}
For every prime $p$ and every $\varepsilon>0$,
\begin{equation}\label{eq:id}
\ZtwoGW(\varepsilon)
 = -\tfrac14\,e^{\varepsilon^2/4}\cosh\!\Bigl(\tfrac{\log p}{2}\Bigr)
   \;-\; \Pp(\varepsilon) \;+\; \Gp(\varepsilon)
   \;-\; \tfrac12(\log\pi)\,A_p(\varepsilon),
\end{equation}
where, with $b_\varepsilon(\delta):=\tfrac{1}{2\varepsilon^2}
-\tfrac{\delta^2}{4\varepsilon^4}$ and $L=\log p$,
\[
\Pp(\varepsilon)
 = \sum_{q,\,m\ge 1}\frac{\log q}{q^{m/2}}\,\frac{\sqrt\pi}{4\pi\varepsilon}
   \Bigl[b_\varepsilon(L-m\log q)\,e^{-(L-m\log q)^2/4\varepsilon^2}
       + b_\varepsilon(L+m\log q)\,e^{-(L+m\log q)^2/4\varepsilon^2}\Bigr],
\]
\[
\Gp(\varepsilon)
 = \frac{1}{2\pi}\int_0^\infty t^2 e^{-\varepsilon^2 t^2}\cos(t L)\,
   \psiR(t)\,dt,
\qquad
A_p(\varepsilon)
 = \frac{1}{2\pi}\int_{-\infty}^\infty t^2 e^{-\varepsilon^2 t^2}\cos(t L)\,dt .
\]
The pole term equals
$-\tfrac18\cdot 2e^{\varepsilon^2/4}\cosh(L/2)$, i.e.\ minus one quarter of
the zeroth-moment pole term, and is $O(\sqrt p\,)$ uniformly for
$\varepsilon\le 1$. The conductor term $-\tfrac12(\log\pi)A_p(\varepsilon)$ comes from
the factor $\pi^{-s/2}$ of the completed zeta function (the half-sum
$\tfrac12\sum_\rho$ carries the prefactor $\tfrac12$); with
$A_p(\varepsilon)=\tfrac{1}{4\sqrt\pi\,\varepsilon^3}
\bigl(1-\tfrac{L^2}{2\varepsilon^2}\bigr)e^{-L^2/4\varepsilon^2}$ it is
$O\!\bigl(\varepsilon^{-5}e^{-L^2/4\varepsilon^2}\bigr)$ as $\varepsilon\downarrow0$;
the conductor contribution $-\tfrac12(\log\pi)A_p$ is numerically
$\approx 8.4\cdot10^{-17}$ for $p=2$ at $\varepsilon=0.05$, negligible on the
scale of the certified margins, but is retained for exactness. In the interval
certificate, which lower-bounds $\mathrm{Cross}=-8\sqrt\pi\,\varepsilon^3\BinfGW$,
the conductor appears as $+(\log\pi)\,h(L^2/4\varepsilon^2)$ with a
\emph{positive} coefficient: the normalisation $8\sqrt\pi\,\varepsilon^3$ turns
$-\tfrac12(\log\pi)A_p$ in $N_{\mathrm{GW}}$ into $-(\log\pi)\,h$, and
$\mathrm{Cross}=-N_{\mathrm{GW}}$ flips this to $+(\log\pi)\,h$, since
$\bigl(1-\tfrac{L^2}{2\varepsilon^2}\bigr)e^{-L^2/4\varepsilon^2}=+h(L^2/4\varepsilon^2)$.
\end{theorem}

\begin{proof}
\emph{Step 1 (zeroth moment).}
The Gaussian test function $h^{(0)}_{p,s}$ is even, entire, and of rapid
decay in horizontal strips, hence admissible
\cite[Thm.~5.12]{IwaniecKowalski04}; for every $s>0$,
\begin{equation}\label{eq:id0}
\begin{aligned}
\tfrac12\sum_\rho h^{(0)}_{p,s}\bigl((\rho-\tfrac12)/i\bigr)
 ={}& e^{s/4}\cosh\!\Bigl(\tfrac{L}{2}\Bigr)
 - \sum_{q,m}\frac{\log q}{q^{m/2}}\,\widehat h^{(0)}_s(m\log q)\\
 &{}+ \frac{1}{2\pi}\int_0^\infty e^{-s t^2}\cos(t L)\,\psiR(t)\,dt
 - \tfrac12(\log\pi)\,A^{(0)}_p(s),
\end{aligned}
\end{equation}
with prime-side transform and conductor integral
\[
\widehat h^{(0)}_s(u)=\tfrac{1}{4\sqrt{\pi s}}
\bigl[e^{-(L-u)^2/4s}+e^{-(L+u)^2/4s}\bigr],
\qquad
A^{(0)}_p(s)=\tfrac{1}{2\pi}\int_{-\infty}^\infty e^{-st^2}\cos(tL)\,dt.
\]
The conductor constant $-\log\pi$ from $\pi^{-s/2}$ does \emph{not} cancel.
Carried through the half-sum $\tfrac12\sum_\rho$ it contributes
$-\tfrac12(\log\pi)A^{(0)}_p(s)$; applying $-\partial_s$ at
$s=\varepsilon^2$ then produces the term $-\tfrac12(\log\pi)A_p(\varepsilon)$
of~\eqref{eq:id}.

\emph{Step 2 (differentiate in $s$).}
Differentiate~\eqref{eq:id0} in $s$ at $s=\varepsilon^2$ and multiply by
$-1$.
On the zero side, termwise differentiation is justified by dominated
convergence, the derivative of each term being bounded by the summable
majorant of Lemma~\ref{lem:conv}; this gives $\ZtwoGW(\varepsilon)$.
On the pole side,
$-\partial_s\,e^{s/4}\cosh(L/2)=-\tfrac14 e^{s/4}\cosh(L/2)$.
On the prime side, for $\delta=L\mp u$,
\[
-\partial_s\bigl[s^{-1/2}e^{-\delta^2/4s}\bigr]
=s^{-1/2}e^{-\delta^2/4s}\bigl(\tfrac{1}{2s}-\tfrac{\delta^2}{4s^2}\bigr),
\]
which yields the bracket $b_\varepsilon$ with prefactor
$\tfrac{1}{4\sqrt\pi\varepsilon}=\tfrac{\sqrt\pi}{4\pi\varepsilon}$;
termwise differentiation is justified directly, the $s$-derivative of the
$(q,m)$ term being bounded for $s\ge s_0>0$ by
$\tfrac{\log q}{q^{m/2}}\,C(s_0)\,(1+(m\log q)^2)\,e^{-(L-m\log q)^2/4s_0}$,
a Gaussian in $m\log q$ that sums absolutely (independently of the later
Theorem~\ref{thm:prime}).
On the archimedean side, differentiation under the integral sign is
justified by the dominating function
$t^2 e^{-s_0 t^2}(\log(2+t)+8)$ for $s\ge s_0>0$ (Lemma~\ref{lem:psi}).
This proves~\eqref{eq:id}.
\end{proof}

\begin{remark}[Numerical sanity check]
At $p\in\{2,5,11,37\}$ and $\varepsilon=0.05$, the two sides of the
zeroth-moment identity~\eqref{eq:id0} agree to $10^{-15}$ and the two sides
of the differentiated identity~\eqref{eq:id} to $10^{-13}$ relative
precision. These checks are not used in the proof.
\end{remark}

% ============================================================
\section{The Archimedean Term}
\label{sec:arch}
% ============================================================

\begin{lemma}[Digamma bounds on the quarter line; proved]
\label{lem:psi}
With $a_n:=n+\tfrac14$ and $\tau:=t/2$,
\[
\psiR(t) = -\gamma_E+\sum_{n\ge 0}
   \Bigl(\tfrac{1}{n+1}-\tfrac{a_n}{a_n^2+\tau^2}\Bigr),
\]
and for all $t\ge 0$:
$|\psiR(t)|\le\log(2+t)+8$ and
$|\psiR'(t)|\le\tfrac{20}{1+t}$.
Moreover $\psiR$ is even, so all odd derivatives vanish at $t=0$, and
$\int_0^\infty S_1(t/2)\,dt=\int_0^\infty t^2 S(t/2)\,dt
=\tfrac{\pi}{2}\psi'(\tfrac14)=\tfrac{\pi}{2}(\pi^2+8G)\approx 27.013$,
where $S_1,S$ are the trigamma/quadrigamma series and $G$ is Catalan's
constant.
\end{lemma}

\begin{proof}
The series follows from
$\psi(z)=-\gamma_E+\sum_{n\ge 0}(\tfrac{1}{n+1}-\tfrac{1}{n+z})$ at
$z=\tfrac14+i\tau$.
\emph{Bound on $\psiR'$.} Differentiating,
$\psiR'(t)=-\tfrac12\Im\psi'(\tfrac14+i\tau)=\tau S_1(\tau)$ with
$S_1(\tau)=\sum_n a_n/(a_n^2+\tau^2)^2$ (since
$\Im\psi'(\tfrac14+i\tau)=-2\tau\sum_n a_n/(a_n^2+\tau^2)^2$; it is
$-\tfrac12\Im\psi'$, \emph{not} $\tfrac12\Re\psi'$, consistent with $\psiR$
even). We bound $\tau S_1(\tau)\le 20/(1+t)$ in two regions.
For $0\le t\le 2$, per-term maximisation
($\max_\tau\tau a_n/(a_n^2+\tau^2)^2=9/(16\sqrt3\,a_n^2)$ at $\tau=a_n/\sqrt3$)
gives $\tau S_1(\tau)\le M_1:=\tfrac{9}{16\sqrt3}\psi'(\tfrac14)<5.59
\le\tfrac{20}{3}\le\tfrac{20}{1+t}$.
For $t\ge2$ (so $\tau\ge1$), comparing the unimodal summand to its integral
plus its peak,
\[
S_1(\tau)\le\int_0^\infty\!\tfrac{a}{(a^2+\tau^2)^2}\,da
+\max_a\tfrac{a}{(a^2+\tau^2)^2}
=\tfrac{1}{2\tau^2}+\tfrac{9}{16\sqrt3\,\tau^3},
\]
so $\tau^2 S_1(\tau)\le\tfrac12+\tfrac{9}{16\sqrt3\,\tau}
\le\tfrac12+\tfrac{9}{16\sqrt3}<0.825$ for $\tau\ge1$; hence
$\psiR'(t)=\tau S_1(\tau)\le 0.825/\tau=1.65/t<\tfrac{20}{1+t}$. This proves the
envelope on all of $[0,\infty)$.

\emph{Bound on $\psiR$.} Write $z=\tfrac14+i\tfrac t2$ and use the recurrence
$\psi(z)=\psi(z+1)-1/z$. Since $\Re(z+1)=\tfrac54\ge\tfrac12$, the standard
digamma estimate $|\psi(w)-\log w|\le 1/|w|$ for $\Re w\ge\tfrac12$
\cite[6.3.18]{AbramowitzStegun} gives
$|\psi(z+1)-\log(z+1)|\le 1/|z+1|\le\tfrac45$. Moreover
$\Re\log(z+1)=\tfrac12\log(\tfrac{25}{16}+\tfrac{t^2}{4})\le\log(2+t)$ (as
$\tfrac{25}{16}+\tfrac{t^2}{4}\le(2+t)^2$), and
$|\Re(1/z)|=\tfrac14/(\tfrac1{16}+\tfrac{t^2}{4})\le 4$. Hence
\[
\begin{aligned}
|\psiR(t)|=|\Re\psi(z)|
&\le|\Re\log(z+1)|+|\psi(z+1)-\log(z+1)|+|\Re(1/z)|\\
&\le\log(2+t)+\tfrac45+4<\log(2+t)+8.
\end{aligned}
\]
Both envelopes are upper bounds on the true functions, so the $I_4$ assembly
(which uses only $|\psiR'|\le20/(1+t)$ and $|\psiR|\le\log(2+t)+8$; the sharp
higher-derivative bounds come from Lemma~\ref{lem:digamma}) remains valid.
The exact integrals reduce, after the substitution $t=2a\lambda$, to
$\tfrac{\pi}{2}\sum_n a_n^{-2}=\tfrac{\pi}{2}\psi'(\tfrac14)$, and
$\psi'(\tfrac14)=\pi^2+8G$.
\end{proof}

For the fourth-derivative bound below we record the exact derivatives of
$\varphi(t):=t^2 e^{-\varepsilon^2 t^2}$,
\begin{align*}
\varphi'   &= (2t-2\varepsilon^2 t^3)\,e^{-\varepsilon^2 t^2}, &
\varphi''  &= (2-10\varepsilon^2 t^2+4\varepsilon^4 t^4)\,e^{-\varepsilon^2 t^2},\\
\varphi''' &= (-24\varepsilon^2 t+36\varepsilon^4 t^3-8\varepsilon^6 t^5)\,e^{-\varepsilon^2 t^2}, &
\varphi''''&= (-24\varepsilon^2+156\varepsilon^4 t^2-112\varepsilon^6 t^4+16\varepsilon^8 t^6)\,e^{-\varepsilon^2 t^2},
\end{align*}
together with the following higher digamma bounds.

\begin{lemma}[Higher digamma bounds; proved]
\label{lem:digamma}
Let $S(\tau):=\sum_{n\ge 0}a_n/(a_n^2+\tau^2)^3$ and
$S_1(\tau):=\sum_{n\ge 0}a_n/(a_n^2+\tau^2)^2$, with $a_n=n+\tfrac14$ and
$\tau=t/2$. Then for all $t\ge 0$,
\[
|\psiR''(t)|\le\tfrac32 S_1(\tau),\qquad
|\psiR'''(t)|\le 3\tau\,S(\tau),\qquad
|\psiR''''(t)|\le\tfrac{15}{2}\,S(\tau),
\]
and moreover
\begin{enumerate}[label=\textup{(\roman*)}]
\item $\displaystyle\int_0^\infty S_1(t/2)\,dt=\int_0^\infty t^2 S(t/2)\,dt
   =\tfrac{\pi}{2}\psi'(\tfrac14)=\tfrac{\pi}{2}(\pi^2+8G)\approx 27.013$;
\item for $\tau\ge 1$, $S(\tau)\le c_S/\tau^4$ with
   $c_S=\tfrac14+\tfrac{125}{216\sqrt5}\le 0.5089$, and
   $S_1(\tau)\le c_{S_1}/\tau^2$ with
   $c_{S_1}=\tfrac12+\tfrac{3\sqrt3}{16}\le 0.8248$;
\item $\displaystyle\sup_{t>0}t^3 S(t/2)\le\sum_{n\ge 0}a_n^{-2}=\psi'(\tfrac14)$;
\item $S(\tau)\le\zeta(5,\tfrac14)$ and $S_1(\tau)\le\zeta(3,\tfrac14)$ for all
   $\tau\ge 0$, both decreasing in $\tau$.
\end{enumerate}
\end{lemma}

\begin{proof}
Termwise differentiation of the series of Lemma~\ref{lem:psi} is justified by
locally uniform convergence; with $f_a(\tau)=a/(a^2+\tau^2)$ and
$\tfrac{d}{dt}=\tfrac12\tfrac{d}{d\tau}$,
\[
f_a''=\frac{-2a(a^2-3\tau^2)}{(a^2+\tau^2)^3},\quad
f_a'''=\frac{24a\tau(a^2-\tau^2)}{(a^2+\tau^2)^4},\quad
f_a''''=\frac{24a(a^4-10a^2\tau^2+5\tau^4)}{(a^2+\tau^2)^5}.
\]
Using $|a^2-3\tau^2|\le 3(a^2+\tau^2)$, $|a^2-\tau^2|\le a^2+\tau^2$, and (with
$a^2=(1-s)u$, $\tau^2=su$) $\max_{0\le s\le 1}|1-12s+16s^2|=5$, so
$|a^4-10a^2\tau^2+5\tau^4|\le 5(a^2+\tau^2)^2$, dividing by $4,8,16$ gives the
three pointwise bounds.
For~(i), substitute $t=2a\lambda$ and use
$\int_0^\infty(1+\lambda^2)^{-2}d\lambda=\pi/4$,
$\int_0^\infty\lambda^2(1+\lambda^2)^{-3}d\lambda=\pi/16$; both reduce to
$\tfrac{\pi}{2}\sum_n a_n^{-2}$.
For~(ii), integral comparison with the maximal summand gives the stated
constants; for~(iii), $\max_{\tau>0}8a\tau^3/(a^2+\tau^2)^3=a^{-2}$ at
$\tau=a$; for~(iv), each summand decreases in $\tau$ and equals the Hurwitz
value at $\tau=0$.
\end{proof}

\begin{lemma}[Forward bound for the $I_4$ residual; proved, certified]\label{lem:resid}
Let $T_1(\varepsilon)=\int_0^\infty|\varphi''''|\,|\psiR|\,dt$ and let
$R_3(\varepsilon)=T_3(\varepsilon)-9\pi\psi'(\tfrac14)$,
$R_4(\varepsilon)=T_4(\varepsilon)-6\pi\psi'(\tfrac14)$ be the remainders of
the two main terms (notation as in Theorem~\ref{thm:arch} below). There are
elementary majorants
\[
T_1(\varepsilon)\le A_1(\varepsilon),\qquad
R_3(\varepsilon)\le A_3(\varepsilon),\qquad
R_4(\varepsilon)\le A_4(\varepsilon),
\]
with endpoint values, and with $A_1$ increasing on $(0,\tfrac14]$ while
$A_3,A_4$ obey $\varepsilon$-uniform bounds (proof),
\[
A_1(\tfrac14)\le 475,\quad A_3(\tfrac14)\le 57,\quad A_4(\tfrac14)\le 14,
\]
so that
\[
A_1(\tfrac14)+A_3(\tfrac14)+A_4(\tfrac14)\le 546\le 1697.1 .
\]
\end{lemma}

\begin{proof}
\emph{Strategy.} We bound $|g''''|$ termwise into three majorants $A_1,A_3,A_4$,
certify their values at $\varepsilon=\tfrac14$ in Arb, and then establish the
$\varepsilon$-uniform bound on $(0,0.05]$ separately (no monotonicity).
With $\varphi(t)=t^2e^{-\varepsilon^2 t^2}$,
$\varphi''''=(-24\varepsilon^2+156\varepsilon^4t^2-112\varepsilon^6t^4
+16\varepsilon^8t^6)e^{-\varepsilon^2t^2}$, so by the triangle inequality and
$|\psiR(t)|\le\log(2+t)+8$ (Lemma~\ref{lem:psi}),
\[
A_1(\varepsilon):=\!\int_0^\infty\!\!\bigl(24\varepsilon^2+156\varepsilon^4t^2
+112\varepsilon^6t^4+16\varepsilon^8t^6\bigr)e^{-\varepsilon^2t^2}
\bigl(\log(2+t)+8\bigr)\,dt
\]
is a majorant for $T_1$; each monomial is a Gaussian moment times a
logarithmic factor, and the interval certificate
(\texttt{cert\_paper8.py}, \texttt{certify\_I4}) integrates this majorant
rigorously in Arb ball arithmetic, giving $A_1(\tfrac14)=474.6\le 475$. For $R_3$, write
$\varphi''=(2-10\varepsilon^2t^2+4\varepsilon^4t^4)e^{-\varepsilon^2t^2}$ and
$|\psiR''|\le\tfrac32 S_1(t/2)$ (Lemma~\ref{lem:digamma}); the constant part
$2$ contributes $9\!\int 2\,S_1(t/2)\,dt=9\pi\psi'(\tfrac14)$, exactly the
subtracted $T_3$ main term, so
\[
R_3(\varepsilon)\le A_3(\varepsilon):=9\!\int_0^\infty\!\!
\bigl(10\varepsilon^2t^2+4\varepsilon^4t^4\bigr)e^{-\varepsilon^2t^2}
S_1(t/2)\,dt,
\]
with $S_1$ the sharp series of Lemma~\ref{lem:digamma} (not the crude
$(1+\tau)^{-2}$ bound); $A_3$ is integrated rigorously in Arb,
$A_3(\tfrac14)=56.1\le 57$. For $R_4$, from
$\varphi'=(2t-2\varepsilon^2t^3)e^{-\varepsilon^2t^2}$ and
$|\psiR'''|\le 3\tau S(\tau)$ the leading $2t$ part contributes
$12\!\int t^2 S(t/2)\,dt=6\pi\psi'(\tfrac14)$ (the $T_4$ main term), leaving
\[
R_4(\varepsilon)\le A_4(\varepsilon):=12\varepsilon^2\!\int_0^\infty\!\!
t^4 e^{-\varepsilon^2t^2}S(t/2)\,dt,
\]
again integrated in Arb, $A_4(\tfrac14)=13.7\le 14$. Summing the three
endpoint values gives $474.6+56.1+13.7=544.4\le1697.1$; the residual bounds
$A_1,A_3,A_4$, the three main terms, and $T_2$ are all reproduced as Arb
outputs by \texttt{certify\_I4} (which also self-tests that the digamma
majorants dominate the true series), so $I_4(\tfrac14)\le3561.1$ is certified,
not merely numerical.

For the certificate range $0<\varepsilon\le 0.05$ we do \emph{not} invoke any
monotonicity of $I_4$; instead the uniform bound is certified directly
(\texttt{certify\_Gw\_uniform}). After the substitution $u=\varepsilon t$,
\[
A_1(\varepsilon)=\varepsilon\!\int_0^\infty\!\!(24+156u^2+112u^4+16u^6)
e^{-u^2}\bigl(\log(2+u/\varepsilon)+8\bigr)\,du,
\]
whose integrand carries the factor $\varepsilon\bigl(\log(2+u/\varepsilon)+8\bigr)$
with $\partial_\varepsilon$-derivative
$\log(2+u/\varepsilon)+8-\tfrac{1}{1+2\varepsilon/u}\ge 7>0$, so $A_1$ is
increasing and $A_1(\varepsilon)\le A_1(0.05)$. Writing $\tau=u/2\varepsilon$,
\[
\begin{aligned}
A_3(\varepsilon)&=18\!\int_0^\infty\!\!(10u+4u^3)e^{-u^2}\,\tau S_1(\tau)\,du
\le 126\,M_1,\\
A_4(\varepsilon)&=96\!\int_0^\infty\!\! u\,e^{-u^2}\,\tau^3 S(\tau)\,du
\le 48\,M_3,
\end{aligned}
\]
where $\int(10u+4u^3)e^{-u^2}=7$, $\int u\,e^{-u^2}=\tfrac12$, and the suprema
are bounded by summing the per-term maxima,
$M_1:=\sup_{\tau\ge0}\tau S_1(\tau)\le\tfrac{9}{16\sqrt3}\,\psi'(\tfrac14)$
(each $\tau\,a_n/(a_n^2+\tau^2)^2$ peaks at $\tau=a_n/\sqrt3$) and
$M_3:=\sup_{\tau\ge0}\tau^3 S(\tau)\le\tfrac18\,\psi'(\tfrac14)$ (peak at
$\tau=a_n$). Hence $I_4(\varepsilon)\le T_2(0.05)+\text{(main)}+A_1(0.05)
+126M_1+48M_3$, giving $G_w(\varepsilon)\le 1586\le 2693$ for all
$\varepsilon\in(0,0.05]$ as a single $\varepsilon$-uniform Arb output. It is
this uniform bound, not monotonicity, that transfers the archimedean constant
to the whole certificate range.
\end{proof}


\begin{theorem}[Explicit archimedean bound; proved]
\label{thm:arch}
For all $0<\varepsilon\le\tfrac14$ and all primes $p\ge 2$ the per-prime bound
$|\Gp(\varepsilon)|\le I_4(\varepsilon)/2\pi(\log p)^4$ holds; on the certificate
range $0<\varepsilon\le0.05$ and at the endpoint $\varepsilon=\tfrac14$ the
certified ceiling $I_4\le 3561.1$ (Lemma~\ref{lem:resid}, no monotonicity) gives
\[
\begin{aligned}
\bigl|\Gp(\varepsilon)\bigr|
 &\;\le\; \frac{I_4(\varepsilon)}{2\pi(\log p)^4}
 \;\le\; \frac{567}{(\log p)^4}
 \;\le\; \frac{1180}{(\log p)^2},\\[2pt]
&I_4(\tfrac14)\le 3561.1,\quad
I_4(\varepsilon)\le 3561.1\ \text{on }(0,0.05]\ (\text{Lemma~\ref{lem:resid}}).
\end{aligned}
\]
In particular $\Gp(\varepsilon)=O(1)$, uniformly in $p$, and the weighted
total obeys
\[
\sum_{p\le 53}(\log p)^2|\Gp(\varepsilon)|\le 2693 .
\]
\end{theorem}

\begin{proof}
Write $\varphi(t)=t^2 e^{-\varepsilon^2 t^2}$ and $g=\varphi\,\psiR$, so
that $2\pi\Gp=\int_0^\infty g(t)\cos(tL)\,dt$.
Since $\varphi(0)=\varphi'(0)=\varphi'''(0)=0$, $\varphi''(0)=2$, and all
odd derivatives of $\psiR$ vanish at $0$ (Lemma~\ref{lem:psi}), one has
$g(0)=g'(0)=0$ and
$g'''(0)=6\,\psiR'(0)=0$.
Four integrations by parts against $\cos(tL)$ therefore leave no boundary
terms and give
\[
\int_0^\infty g(t)\cos(tL)\,dt
 = \frac{1}{L^4}\int_0^\infty g''''(t)\cos(tL)\,dt,
\qquad
|\Gp|\le\frac{I_4(\varepsilon)}{2\pi L^4},
\quad
I_4(\varepsilon):=\int_0^\infty|g''''(t)|\,dt.
\]
\emph{Assembly of $I_4$.} By Leibniz,
\[
g''''=\varphi''''\psiR+4\varphi'''\psiR'+6\varphi''\psiR''
      +4\varphi'\psiR'''+\varphi\,\psiR'''',
\]
and we bound the five integrals $T_1,\dots,T_5$ separately, writing
$E:=e^{-\varepsilon^2 t^2}$ and
$M_j:=\int_0^\infty t^j E\,dt=\Gamma(\tfrac{j+1}{2})/(2\varepsilon^{j+1})$.
\begin{itemize}[leftmargin=2em]
\item $T_1=\int|\varphi''''|\,|\psiR|\,dt$: with $|\psiR|\le\log(2+t)+8$
  (Lemma~\ref{lem:psi}) and $\log(2+t)\le\log(2+1/\varepsilon)+\log(1+\varepsilon t)$,
  each of the four monomials of $\varphi''''$ contributes a term carrying a
  prefactor $\varepsilon^{2k+2}$, all $O(\varepsilon\log(1/\varepsilon))$.
\item $T_2=4\int|\varphi'''|\tfrac{20}{1+t}\,dt
  \le 80(24\varepsilon^2 M_0+36\varepsilon^4 M_2+8\varepsilon^6 M_4)=O(\varepsilon)$,
  using $t^j/(1+t)\le t^{j-1}$.
\item $T_3=6\int|\varphi''|\tfrac32 S_1(t/2)\,dt
  \le 9\bigl[2\!\int_0^\infty\! S_1+10\varepsilon^2\!\int t^2 E S_1
  +4\varepsilon^4\!\int t^4 E S_1\bigr]$; the first integral is exact by
  Lemma~\ref{lem:digamma}(i), the remainders use the split at $t=2$ with
  Lemma~\ref{lem:digamma}(ii),(iv) and are $O(\varepsilon)$. Main term
  $18\cdot\tfrac{\pi}{2}\psi'(\tfrac14)=9\pi\psi'(\tfrac14)\approx 486.2$.
\item $T_4=4\int|\varphi'|\tfrac{3t}{2}S(t/2)\,dt
  \le 12\int t^2 E\,S(t/2)\,dt+O(\varepsilon^0)$; on $[0,2]$ use
  $t^5 S(t/2)\le t^2\psi'(\tfrac14)$ (Lemma~\ref{lem:digamma}(iii)) and on
  $[2,\infty)$ the decay $t^5 S(t/2)\le 16 c_S t$. Main term
  $12\cdot\tfrac{\pi}{2}\psi'(\tfrac14)=6\pi\psi'(\tfrac14)\approx 324.2$,
  remainder $\le\tfrac83\psi'(\tfrac14)\cdot 12\varepsilon^2+96 c_S\le 34.4+48.9$
  on $(0,\tfrac14]$.
\item $T_5=\int\varphi\,\tfrac{15}{2}S(t/2)\,dt
  \le\tfrac{15}{2}\int t^2 S(t/2)\,dt=\tfrac{15\pi}{4}\psi'(\tfrac14)\approx 202.6$
  (using $E\le 1$).
\end{itemize}
Evaluated at $\varepsilon=\tfrac14$ and rounded up, the five rows sum to the
conservative endpoint ceiling $I_4(\tfrac14)\le 3561.1$, which the proof uses.
The accompanying Arb certificate \texttt{certify\_I4} recomputes the same
archimedean budget and outputs the sharper endpoint enclosure
$I_4(\tfrac14)\le 2408.17$ (together with the $\varepsilon$-uniform bound on
$(0,0.05]$ of Lemma~\ref{lem:resid}, no monotonicity, used in the weighted
total); this sharper enclosure is not needed for the proof, it verifies the
conservative ceiling used below.
The residual row is
bounded \emph{forward} by Lemma~\ref{lem:resid}
($A_1+A_3+A_4\le 544.4$ at $\varepsilon=\tfrac14$), well within the $1697.1$
budget, so the table summarizes a proof already present in the bounds rather
than introducing $1697.1$ by subtraction:
\begin{center}
\renewcommand{\arraystretch}{1.2}
\begin{tabular}{l|l|l}
term & majorant at $\varepsilon=\tfrac14$ & upper bound \\
\hline
$T_3$ main $=9\pi\psi'(\tfrac14)$ & $9\pi(\pi^2+8G)$ & $\le 486.3$ \\
$T_4$ main $=6\pi\psi'(\tfrac14)$ & $6\pi(\pi^2+8G)$ & $\le 324.2$ \\
$T_5=\tfrac{15\pi}{4}\psi'(\tfrac14)$ & $\tfrac{15\pi}{4}(\pi^2+8G)$ & $\le 202.7$ \\
$T_2$ & $80(24\varepsilon^2 M_0+36\varepsilon^4 M_2+8\varepsilon^6 M_4)$ & $\le 850.8$ \\
$T_1$ and the $T_3,T_4$ remainders & forward via Lem.~\ref{lem:resid} ($\le 544.4$) & $\le 1697.1$ (budget) \\
\hline
$I_4(\tfrac14)$ & total & $\le 3561.1$ \\
\end{tabular}
\end{center}
Hence $567/L^4\le 567/(L^2(\log 2)^2)\le 1180/L^2$ as a per-prime bound.
For the weighted total the rounded constant $567=\lceil I_4(\tfrac14)/2\pi\rceil$
must be avoided (it would give $567\cdot 4.7514=2694.0>2693$); using the
unrounded value $I_4(\tfrac14)/2\pi=3561.1/2\pi=566.77$,
\[
\sum_{p\le 53}(\log p)^2|\Gp(\varepsilon)|
 \;\le\; \frac{I_4(\tfrac14)}{2\pi}\sum_{p\le 53}(\log p)^{-2}
 \;=\; \frac{3561.1}{2\pi}\cdot 4.7514 \;<\; 2693 .
\]
\end{proof}

\begin{remark}
The bound is conservative by roughly two orders of magnitude. Numerically,
the weighted terms satisfy $(\log p)^2\Gp(\varepsilon)\in[0.26,0.76]$ for
$p\le 53$ and $\varepsilon\in[0.01,0.1]$, and the residuals numerically
appear to remain bounded, consistent with finite $O(1)$ limiting behaviour
as $\varepsilon\to 0$.
The point is that $\Gp$ is subdominant by three powers of $\varepsilon$
against the prime-side main term of order $\varepsilon^{-3}$.
\end{remark}

% ============================================================
\section{The Prime Side}
\label{sec:prime}
% ============================================================

\begin{theorem}[Eigenterm extraction and cross-prime control; proved]
\label{thm:prime}
Fix $\kappa\ge 2$ and a prime $p\le\kappa$, and let
\[
\delta_p := \min\bigl\{\,|L-m\log q| : q\text{ prime},\ m\ge 1,\
   (q,m)\ne(p,1)\,\bigr\}>0,
\qquad L=\log p.
\]
Then for all $0<\varepsilon\le\tfrac14$,
\[
\Pp(\varepsilon)
 = \frac{\log p}{\sqrt p\;8\sqrt\pi\,\varepsilon^3} \;+\; E_p(\varepsilon),
\qquad
|E_p(\varepsilon)|
 \le \frac{c_p}{\varepsilon^5}\,e^{-\delta_p^2/4\varepsilon^2}
   + \frac{c_p^{\mathrm{far}}}{\varepsilon^3}\,e^{-1/16\varepsilon^2},
\]
with explicit constants
$c_p=\tfrac{3\sqrt\pi}{16\pi}
\sum_{(q,m)\in\mathcal N_p}\tfrac{\log q}{q^{m/2}}$ (the direct near sum) and
$c_p^{\mathrm{far}}=\tfrac{5\sqrt\pi}{16\pi}
(1.506\,p+1.04\sqrt p+0.76)$, where
$\mathcal N_p=\{(q,m)\ne(p,1):|L-m\log q|<1\}$ and $n_p=\#\mathcal N_p$.
\end{theorem}

\begin{proof}
\emph{Eigenterm.}
The term $(q,m)=(p,1)$ has $\delta_-=0$, hence
$b_\varepsilon(0)=\tfrac{1}{2\varepsilon^2}$ and contributes exactly
\[
\frac{\log p}{\sqrt p}\cdot\frac{\sqrt\pi}{4\pi\varepsilon}
\cdot\frac{1}{2\varepsilon^2}
=\frac{\log p}{\sqrt p\,8\sqrt\pi\,\varepsilon^3},
\]
with no multiplicative error; its $\delta_+=2L$ companion is absorbed below.

\emph{Near zone.}
For $(q,m)\in\mathcal N_p$ one has $\delta_p\le|\delta_-|<1$, so for
$\varepsilon\le 1$, $|b_\varepsilon(\delta_-)|\le\tfrac{3}{4\varepsilon^4}$
and $e^{-\delta_-^2/4\varepsilon^2}\le e^{-\delta_p^2/4\varepsilon^2}$;
summing the direct near sum $\sum_{(q,m)\in\mathcal N_p}\log q/q^{m/2}$
(not $n_p\max$) gives the first bound.

\emph{Far zone.}
For any term with $\delta:=|L\mp m\log q|\ge 1$,
$|b_\varepsilon(\delta)|e^{-\delta^2/4\varepsilon^2}
\le\tfrac{5}{4\varepsilon^2}e^{-\delta^2/16\varepsilon^2}e^{-1/16\varepsilon^2}$,
and for $\varepsilon\le\tfrac14$, $e^{-\delta^2/16\varepsilon^2}\le
e^{-\delta}$.
Writing $n=q^m$, the three far families are summed against $e^{-\delta}$:
for $n>ep$, $\sum_n\Lambda(n)n^{-3/2}p\le 1.506\,p$ (the full Dirichlet
series $\sum_n\Lambda(n)n^{-3/2}=-\zeta'(\tfrac32)/\zeta(\tfrac32)=1.5052\ldots\le 1.506$
is a majorant); for $n<p/e$,
$\sqrt p\,\psi(p)/p\le 1.04\sqrt p$ by Chebyshev~\cite{RosserSchoenfeld62};
for the $\delta_+$ terms, $\tfrac1p\sum_n\Lambda(n)n^{-3/2}\le 0.76$.
Multiplying by $\tfrac{\sqrt\pi}{4\pi\varepsilon}$ yields
$c_p^{\mathrm{far}}$.
\end{proof}

\begin{remark}
For $\kappa=53$ the binding proximity is
$\dmin(53)=\min_{p\le 53}\delta_p=0.031749$, attained by the pair
$31\leftrightarrow 2^5=32$; the weighted near-zone and far-zone constants are
\[
\sum_{p\le 53}(\log p)^2 c_p=163.6\ (=\lceil163.611\rceil,\ \text{direct near sum}),
\qquad
\sum_{p\le 53}(\log p)^2 c_p^{\mathrm{far}}\le 1385.
\]
\end{remark}

% ============================================================
\section{Negativity at the Reference Parameter}
\label{sec:ref}
% ============================================================

Summing the identity of Theorem~\ref{thm:id} against the weights
$(\log p)^2$ over $p\le\kappa$ gives
\begin{equation}\label{eq:assembly}
\BinfGW(\kappa,\varepsilon)
 = -\,M(\varepsilon) + \widetilde R(\varepsilon),
\qquad
M(\varepsilon)=\frac{S_3(\kappa)}{8\sqrt\pi\,\varepsilon^3},
\quad
|\widetilde R(\varepsilon)|\le R(\varepsilon),
\end{equation}
where $-M(\varepsilon)$ collects the diagonal eigenterms of
Theorem~\ref{thm:prime} (each contributing
$-(\log p)^2\cdot\tfrac{\log p}{\sqrt p\,8\sqrt\pi\varepsilon^3}$, with the
overall sign from $\BinfGW=\sum(\log p)^2(\,\cdots-\Pp+\Gp)$), and
$R(\varepsilon)$ bounds, respectively, the weighted near- and far-zone
cross-prime errors of Theorem~\ref{thm:prime}, the weighted archimedean
total of Theorem~\ref{thm:arch}, and the weighted pole term:
\[
R(\varepsilon)
 = \frac{C_{\mathrm{near}}(\kappa)}{\varepsilon^5}\,e^{-\dmin^2/4\varepsilon^2}
 + \frac{C_{\mathrm{far}}(\kappa)}{\varepsilon^3}\,e^{-1/16\varepsilon^2}
 + G_w(\kappa)
 + \frac{e^{1/4}}{4}\sum_{p\le\kappa}(\log p)^2\cosh\tfrac{\log p}{2},
\]
with $C_{\mathrm{near}}(\kappa)=\sum_{p\le\kappa}(\log p)^2 c_p
=\tfrac{3\sqrt\pi}{16\pi}W_{\mathrm{near}}(\kappa)$ (the direct near sum of
Lemma~\ref{lem:assembly}(b), with no $n_p\max$ inflation),
$C_{\mathrm{far}}(\kappa)=\sum_{p\le\kappa}(\log p)^2 c_p^{\mathrm{far}}$, and
\[
G_w(\kappa)=\sup_{0<\varepsilon\le 0.05}\;\sum_{p\le\kappa}(\log p)^2|\Gp(\varepsilon)|
\]
the (uniform, $\varepsilon$-independent) weighted
archimedean total bounded by Theorem~\ref{thm:arch} on the certified range
$(0,0.05]$ ($\le 2693$ for
$\kappa=53$; the supremum is taken over the range where $I_4\le3561.1$ is
certified, not over $(0,\tfrac14]$). For $\kappa=53$, $C_{\mathrm{near}}(53)=163.6$ (direct sum) and
$C_{\mathrm{far}}(53)\le 1385$.
Hence $\BinfGW<0$ whenever $R(\varepsilon)<M(\varepsilon)$.
We state the closed-form and certified ranges separately.

\subsection{Closed-form core}

\begin{theorem}[Closed-form negativity threshold; proved]
\label{thm:closedform}
For $\kappa=53$, $\BinfGW(53,\varepsilon)<0$ for all
$0<\varepsilon\le\varepsilon_0(53):=0.004$.
More generally, for every $\kappa\ge 2$ and every
$0<\varepsilon\le\min\{0.05,\dmin/2\}$ with $R(\varepsilon)<M(\varepsilon)$,
and for all smaller $\varepsilon$, one has $\BinfGW(\kappa,\varepsilon)<0$.
\end{theorem}

\begin{proof}
By~\eqref{eq:assembly} it suffices that $R(\varepsilon)<M(\varepsilon)$, and
that the inequality propagates downward in $\varepsilon$.
Each error-to-main ratio is non-decreasing on the stated range: the
near-zone ratio is proportional to
$\varepsilon^{-2}e^{-\dmin^2/4\varepsilon^2}$ (increasing for
$\varepsilon<\dmin/2$), the far-zone ratio to $e^{-1/16\varepsilon^2}$, the
archimedean ratio to $\varepsilon^3$ (since $G_w(\varepsilon)\le 2693$
holds uniformly, Lemma~\ref{lem:resid}, while $M\propto\varepsilon^{-3}$), and the pole ratio to
$\varepsilon^3 e^{\varepsilon^2/4}$; all increasing.
Hence $R(\varepsilon_0)<M(\varepsilon_0)$ at some admissible $\varepsilon_0$
implies the same for all $0<\varepsilon\le\varepsilon_0$.
For $\kappa=53$: $S_3(53)=87.4247$, $\pi(53)=16$, $\dmin=0.031749$, and the
weighted constants of \S\,\ref{sec:prime}, \S\,\ref{sec:arch}.
At $\varepsilon_0=0.004\le\dmin/2$ one computes $M=9.63\cdot10^{7}$,
near-zone $\le 3.6\cdot10^{7}$, far-zone $\le 10^{-1000}$, archimedean
$G_w(53)\le 2693$, pole $\le 102$, so $R<M$ with margin, and the
monotonicity step extends the conclusion to all $0<\varepsilon\le 0.004$.
\end{proof}

\subsection{Certified extension to the reference width}

The closed-form criterion is deliberately conservative: it bounds every
near-zone term by the global minimum $\dmin$ and uses the crude archimedean
constant.
Replacing the closed-form cross-prime bound by the exact (finitely
truncated, with rigorously negligible tail) absolute cross sums, while
keeping the proven archimedean constant of Theorem~\ref{thm:arch} and the
pole bound, yields the following certified statement.

\begin{theorem}[Negativity at the reference parameter; proved, certified]
\label{thm:ref}
$\BinfGW(53,\varepsilon)<0$ for all $0<\varepsilon\le 0.05$.
The Arb interval criterion holds on $[10^{-9},0.047]$ with margin $0.83$ and on
$[0.047,0.05]$ with margin $24.29$; the remaining $(0,10^{-9})$ is covered by
Theorem~\ref{thm:closedform}, so together these prove $(0,0.05]$. It uses exact finite sums over all prime
powers with $|\log n\mp\log p|<3$ in 256-bit ($\approx$77-digit) Arb interval arithmetic. The
discarded prime powers are bounded by a Gauss--Chebyshev majorant: grouping
them by unit $\log n$-bins, the bin count is at most
$\psi(e^{k+1})\le1.04\,e^{k+1}$ (Chebyshev), $\Lambda(n)\le\log n<k+1$, and
$|h|\le(2x-1)e^{-x}$ with $x=(\log n\mp\log p)^2/4\varepsilon^2$; the quadratic
Gaussian decay $e^{-x}$ dominates the single-exponential bin growth, so the
tail series is super-exponentially convergent and rigorously below
$10^{-380}$.
\end{theorem}

\begin{proof}
Multiply Theorem~\ref{thm:id} by $8\sqrt\pi\varepsilon^3(\log p)^2$ and sum,
so the prime side becomes $S_3(53)+\mathrm{Cross}(\varepsilon)$ with the
cross terms expressed through $h(x):=(1-2x)e^{-x}$ evaluated at
$x_{n,\pm}=(\log n\mp\log p)^2/4\varepsilon^2$.
On any interval $[a,b]$ the argument of each $x$ ranges in $[x(b),x(a)]$.
In the normalised lower-bound expression for
$\mathrm{Cross}=-8\sqrt\pi\,\varepsilon^3\BinfGW$, both the prime cross terms and
the conductor term $+(\log\pi)\,h(x_{\ell_p})$ have positive coefficients; each
is therefore bounded below by its coefficient times
$\min_{x\in[x(b),x(a)]}h(x)$, giving a
closed-form interval lower bound
$\mathrm{LB}[a,b]$.
With the pole term $-\tfrac14 e^{\varepsilon^2/4}\cosh(\log p/2)<0$ and the
archimedean side controlled by Theorem~\ref{thm:arch} (weighted total
$\le 2693$), the interval criterion
$S_3(53)+\mathrm{LB}[a,b]>8\sqrt\pi b^3\cdot 2693$ implies
$\BinfGW(53,\varepsilon)<0$ on $[a,b]\cap(0,\tfrac14]$.
Here the archimedean constant $2693$ is the weighted total
$\bigl\lceil \tfrac{I_4(1/4)}{2\pi}\sum_{p\le53}(\log p)^{-2}\bigr\rceil
=\lceil \tfrac{3561.1}{2\pi}\cdot4.7514\rceil$, with $I_4(\tfrac14)\le3561.1$
established through the explicit integration-by-parts bounds of
Lemma~\ref{lem:resid}; this bound is \emph{recertified in Arb ball arithmetic}
by the interval certificate (\texttt{certify\_I4}), which outputs
$I_4(\tfrac14)\le2408.17$ and hence $G_w\le2693$, so $2693$ is a certified
output, not a numerical input.
Evaluating the exact finite sums at 256-bit precision gives the two stated
single-interval certificates with margins $0.83$ and $24.29$.
\end{proof}

\begin{remark}[Certificate data]
The certified inequalities of Theorem~\ref{thm:ref} were produced by the
interval certificate in the repository~\cite{Code}, run in Arb ball
arithmetic (FLINT~3, Python bindings) with directed outward (midpoint--radius)
rounding on all operations at 256-bit working precision ($\approx$77 digits). All interval endpoints are
decimal rationals. The two reference intervals cover $[10^{-9},0.047]$ (margin
$0.83$) and $[0.047,0.05]$ (margin $24.29$), i.e.\ the interval certificate
covers $[10^{-9},0.05]$; the remaining $(0,10^{-9})$ is covered by the
closed-form Theorem~\ref{thm:closedform} (valid for $\varepsilon\le0.004$), so
the \emph{union} of the two arguments gives $(0,0.05]$. The extension to
$\varepsilon\le 0.0819$
uses $378$ chained intervals of the same form. In each interval the finite
sum runs over all prime powers with $|\log n\mp\log p|<3$, the discarded
tail being rigorously below $10^{-380}$. The per-interval lower bound, right
side, and margin are tabulated in the script output, archived under the
release tag of~\cite{Code}.
The certified run used $\texttt{PREC\_BITS}=256$, $\kappa=53$,
$\texttt{GW\_ARCH\_BOUND}=2693$, $\texttt{TAIL\_CUTOFF}=3$, covering intervals
$[10^{-9},0.047]$ and $[0.047,0.05]$, parameter-hash \texttt{2b405227},
source-hash \texttt{d080b67e} (the self-hash of the certificate script,
regenerated at the release tag). The certified margins are $0.8345$ and
$24.2883$.
\end{remark}

\begin{corollary}[Full-object curvature positivity; proved]
\label{cor:curv}
The full Guinand--Weil second-moment curvature is positive,
$-2\BinfGW(53,\varepsilon)>0$ for all $0<\varepsilon\le 0.05$.
The operator curvature $O''(\tfrac12)=-2\Binfline$ is obtained from this once
the off-line correction $\Eoff$ is bounded; see
Corollary~\ref{cor:curvline} after Theorem~\ref{thm:offline}.
\end{corollary}

\begin{remark}[Numerical illustration at the reference parameter]
The proved sign $\BinfGW<0$ is made concrete through the on-line proxy at
$(\kappa,\varepsilon,N)=(53,0.05,100)$, purely as an illustration; the
proof above does not rely on this evaluation. The prime side is dominated
by the diagonal eigenterms $\log p/(\sqrt p\,8\sqrt\pi\,\varepsilon^3)$,
summing to $S_3(53)=87.4247$ before cross-prime corrections, the tightest
of which sits at the near-zone gap $\dmin(53)=0.031749$ between $\log 32$
and $\log 31$ (the prime power $2^5$ against the prime $31$). Each pole term
$-\tfrac14 e^{\varepsilon^2/4}\cosh(\tfrac12\log p)$ is negative, and the
archimedean budget is bounded by $2693$ (with true weighted value
$\approx 5.0$). The net second-moment curvature sum evaluates to
$B(53,0.05,100)=-19\,342.5$, hence $O''(\tfrac12)=+38\,685.1$, in
agreement with the certified sign. The illustration uses the first $N=100$
ordinates; the theorems themselves are parameter-symbolic and do not depend
on a finite zero list.
\end{remark}

\begin{remark}[Certified continuation]
Chaining $378$ analytic interval certificates of the same type extends the
certified range to $\varepsilon\le 0.0819$, where the exact margin crosses
zero against the present constants.
Above $\varepsilon\approx 0.07$ the binding term is the archimedean budget
(true weighted value $\approx 5.0$ versus the bound $2693$); a sharper
archimedean constant would extend certification further.
The certified sign $\BinfGW(53,\varepsilon)<0$ holds throughout, and the
on-line proxy value is negative numerically for all tested
$\varepsilon\in[0.005,0.60]$, with $B(53,0.05,100)=-19\,342.5$.
\end{remark}

% ============================================================
\section{Uniformity in the Cutoff}
\label{sec:uniform}
% ============================================================

For general $\kappa$ let $\dmin(\kappa)$ denote the minimal distance
$|\log p-m\log q|$ over primes $p\le\kappa$ and prime powers $q^m\ne p$
(with no restriction $q^m\le\kappa$).

\begin{lemma}[Elementary gap bound; proved]
\label{lem:gap}
For all $\kappa\ge 2$, $\dmin(\kappa)\ge\dfrac{1}{e\kappa}$.
\end{lemma}

\begin{proof}
Only pairs with $|\log p-m\log q|<1$ matter, i.e.\
$q^m\in(p/e,ep)\subset(0,e\kappa)$.
Since $p$ is prime and $q^m\ne p$, the integers $p$ and $q^m$ are distinct,
so $|p-q^m|\ge 1$; the mean value theorem for $\log$ gives
$|\log p-m\log q|=|p-q^m|/\xi\ge 1/(e\kappa)$ with $\xi<e\kappa$.
\end{proof}

\begin{lemma}[Global constant assembly; proved]\label{lem:assembly}
With the Chebyshev-type inputs $\vartheta(x)>0.84x$ $(x\ge101)$,
$\psi(x)\le1.04x$ and $\pi(x)<1.25506\,x/\log x$
of~\cite{RosserSchoenfeld62}, for $\kappa\ge202$:
\begin{enumerate}[label=\textup{(\alph*)}]
\item $S_3(\kappa)\ge 0.32\,\sqrt\kappa\,(\log\tfrac\kappa2)^2$;
\item $W_{\mathrm{near}}(\kappa):=\sum_{p\le\kappa}(\log p)^2
      \sum_{(q,m)\in\mathcal N_p}\tfrac{\log q}{q^{m/2}}
      \le 5.85\,\kappa^{3/2}\log\kappa$
      (the direct near-zone sum, \emph{not} $\sum_p n_p\max$);
\item $G_w(\kappa)\le 1481\,\kappa/\log\kappa$;
\item $|B_{\mathrm{pol}}(\kappa,\varepsilon)|\le 0.32\,\kappa^{3/2}\log\kappa$
      for $\varepsilon\le\tfrac14$.
\end{enumerate}
Moreover each ratio of the four error classes to $S_3(\kappa)/4$ is strictly
decreasing on $[202,\infty)$.
\end{lemma}

\begin{proof}
\emph{(a)} Restricting to $\kappa/2<p\le\kappa$ and using
$\sqrt p>\sqrt{\kappa/2}$, $\log p>\log(\kappa/2)$,
\[
S_3(\kappa)\ge(\log\tfrac\kappa2)^2\,\kappa^{-1/2}\bigl(\vartheta(\kappa)-\vartheta(\kappa/2)\bigr)
\ge(\log\tfrac\kappa2)^2\kappa^{-1/2}\bigl(0.84\kappa-1.04\cdot\tfrac\kappa2\bigr)
=0.32\sqrt\kappa(\log\tfrac\kappa2)^2,
\]
since $0.84-0.52=0.32$ and $\kappa/2\ge101$.
\emph{(b)} For fixed $p$ the near zone is the window $|L-m\log q|<1$ of
Theorem~\ref{thm:prime}, i.e.\ every near prime power lies in $(p/e,ep)$, so
the \emph{direct} sum $\sum_{(q,m)\in\mathcal N_p}\tfrac{\log q}{q^{m/2}}
=\sum\Lambda(n)/\sqrt n\le e^{1/2}p^{-1/2}\psi(pe)\le1.04\,e^{3/2}\sqrt p$
(bounding the sum itself, with no $n_p\max$ inflation);
then $\sum_{p\le\kappa}(\log p)^2\sqrt p\le\sqrt\kappa\,(\log\kappa)^2\,\pi(\kappa)
\le\sqrt\kappa\,(\log\kappa)^2\cdot1.25506\,\kappa/\log\kappa
=1.25506\,\kappa^{3/2}\log\kappa$,
and $1.04\,e^{3/2}\cdot1.25506=5.849\le5.85$, so $W_{\mathrm{near}}\le5.85\kappa^{3/2}\log\kappa$.
\emph{(c)} $\sum_{p\le\kappa}\ell_p^{-2}\le\pi(\kappa)/(\log2)^2$ with
$\Gp=O(1)$ from Theorem~\ref{thm:arch}: the per-prime archimedean weight is
$\le I_4(\tfrac14)/2\pi\le 566.77$, and
$566.77\cdot1.25506/(\log2)^2\le1481$.
\emph{(d)} $\cosh(\ell_p/2)\le\sqrt p$ and
$\tfrac14e^{1/64}\cdot1.25506=0.3187\le0.32$.
\emph{Monotonicity.} Writing $u=\log\kappa$, $v=\log(\kappa/2)=u-\log2$, the
four ratios are $R_{\mathrm{near}}\!\propto u/(\kappa v^2)$,
$R_{\mathrm{arch}}\!\propto\kappa^{-5/2}u^{-5/2}v^{-2}$,
$R_{\mathrm{pole}}\!\propto\kappa^{-2}u^{-1/2}v^{-2}$, and
$R_{\mathrm{far}}$ carries the factor $e^{-e^2\kappa^2u/2}$.
Their logarithmic derivatives are
\[
\tfrac{d}{d\kappa}\log R_{\mathrm{near}}=\tfrac1{\kappa u}-\tfrac1\kappa-\tfrac2{\kappa v}
=-0.0062<0,\quad
\tfrac{d}{d\kappa}\log R_{\mathrm{arch}}=-\tfrac2\kappa-\tfrac5{2\kappa u}-\tfrac2{\kappa v}
=-0.0144<0,
\]
\[
\tfrac{d}{d\kappa}\log R_{\mathrm{pole}}=-\tfrac2\kappa-\tfrac1{2\kappa u}-\tfrac2{\kappa v}
=-0.0125<0\quad(\text{all at }\kappa=202),
\]
each manifestly negative for all $\kappa\ge202$ since every summand after the
first is negative and the first, $1/(\kappa\log\kappa)$, is dominated by
$1/\kappa$; $R_{\mathrm{far}}$ decreases trivially. Hence all four ratios
decrease on $[202,\infty)$.
\end{proof}


\begin{theorem}[Uniform negativity threshold; proved]
\label{thm:uniform}
Let $\varepsilon_0(\kappa):=\dfrac{1}{4e\kappa\sqrt{\log\kappa}}$.
Then for every $\kappa\ge 202$ and every $0<\varepsilon\le\varepsilon_0(\kappa)$,
\[
\BinfGW(\kappa,\varepsilon)\;<\;0 .
\]
\end{theorem}

\begin{proof}
Fix $\kappa\ge 202$ and $\varepsilon\le\varepsilon_0:=\varepsilon_0(\kappa)$.
We show each of the four error classes in $R(\varepsilon)$ is below
$S_3(\kappa)/4$ at $\varepsilon_0$; since every bound below is increasing in
$\varepsilon$, this covers the whole interval.
By Lemma~\ref{lem:assembly},
$S_3(\kappa)\ge 0.32\sqrt\kappa(\log\tfrac\kappa2)^2$ and the assembled
weighted constants
$W_{\mathrm{near}}(\kappa)\le 5.85\,\kappa^{3/2}\log\kappa$,
$G_w(\kappa)\le 1481\,\kappa/\log\kappa$,
$|B_{\mathrm{pol}}(\kappa,\varepsilon)|\le 0.32\,\kappa^{3/2}\log\kappa$.

\emph{Near zone.}
By Lemma~\ref{lem:gap}, every near pair has
$x=\delta^2/4\varepsilon_0^2\ge\dmin^2/4\varepsilon_0^2\ge 4\log\kappa\ge\tfrac12$,
so $(1+2x)e^{-x}\le 2.95\,\kappa^{-2}$, whence
$\mathrm{Near}\le 17.3\,\kappa^{-1/2}\log\kappa$ (the value $6.46$ at
$\kappa=202$).

\emph{Far zone.}
$\mathrm{Far}\le 20.5\,\kappa^{3/2}\log\kappa\;e^{-e^2\kappa^2\log\kappa/2}$
with $e^{-1/(32\varepsilon_0^2)}=e^{-e^2\kappa^2\log\kappa/2}$; at $\kappa=202$
this is below $10^{-100}$ and stays $<0.08\sqrt\kappa(\log\tfrac\kappa2)^2$
for all $\kappa\ge 2$.

\emph{Archimedean.}
$8\sqrt\pi\varepsilon_0^3 G_w(\kappa)\le 16.4/(\kappa^2(\log\kappa)^{5/2})$,
which equals $6.2\cdot10^{-6}$ at $\kappa=202$ and is decreasing in $\kappa$.

\emph{Pole.}
$8\sqrt\pi\varepsilon_0^3|B_{\mathrm{pol}}|
\le 3.53\cdot10^{-3}\,\kappa^{-3/2}(\log\kappa)^{-1/2}$, which equals
$5.4\cdot10^{-7}$ at $\kappa=202$ and is decreasing in $\kappa$.

At the endpoint $\kappa=202$ each of the four error classes is below
$S_3(\kappa)/4\ge 0.08\sqrt\kappa(\log\tfrac\kappa2)^2=24.2$ (near $6.46$,
far $<10^{-100}$, archimedean $6.2\cdot10^{-6}$, pole $5.4\cdot10^{-7}$), so
their sum is below $S_3(202)$. The ratio of each class to $S_3(\kappa)/4$ is
decreasing for $\kappa\ge 202$, so the comparison persists:
\[
\frac{\mathrm{Near}}{S_3/4}\le\frac{216.25\,\log\kappa}{\kappa(\log\tfrac\kappa2)^2}
\quad(=0.2668\text{ at }\kappa=202),
\qquad
\frac{\mathrm{Far}}{S_3/4}\le \frac{257\,\kappa\log\kappa}{(\log\tfrac\kappa2)^2}
e^{-e^2\kappa^2\log\kappa/2},
\]
\[
\frac{8\sqrt\pi\varepsilon_0^3 G_w}{S_3/4}
 \le\frac{205}{\kappa^{5/2}(\log\kappa)^{5/2}(\log\tfrac\kappa2)^2},
\qquad
\frac{8\sqrt\pi\varepsilon_0^3|B_{\mathrm{pol}}|}{S_3/4}
 \le\frac{0.045}{\kappa^{2}(\log\kappa)^{1/2}(\log\tfrac\kappa2)^2},
\]
each ratio decreasing on $[202,\infty)$ by the logarithmic-derivative
computation of Lemma~\ref{lem:assembly}.

The monotonicity in $\varepsilon$ and the condition $x\ge\tfrac12$ hold on
all of $(0,\varepsilon_0]$ because $\varepsilon_0\le\dmin/\sqrt2$; hence
$-8\sqrt\pi\varepsilon^3\BinfGW\ge S_3-\mathrm{Near}-\mathrm{Far}
-8\sqrt\pi\varepsilon^3(G_w+|B_{\mathrm{pol}}|)>0$ throughout.
\end{proof}

\begin{corollary}[Uniform full-object curvature positivity; proved]
\label{cor:curvuniform}
For every $\kappa\ge 202$ and $0<\varepsilon\le\varepsilon_0(\kappa)$, the
full Guinand--Weil curvature is positive, $-2\BinfGW(\kappa,\varepsilon)>0$.
The operator curvature $O''(\tfrac12)=-2\Binfline$ follows in the certified
range $\varepsilon_{\mathrm{off}}\le\varepsilon\le\varepsilon_0(\kappa)$ by
the off-line transfer of Theorem~\ref{thm:offline}
(Corollary~\ref{cor:curvline}).
\end{corollary}

\begin{remark}[Calibration]
The threshold $\varepsilon_0(\kappa)$ is conservative by design.
Numerically $\kappa\,\dmin(\kappa)\in[1.6,2.5]$ over $\kappa\in[53,3203]$,
so Lemma~\ref{lem:gap} is sharp up to the constant; at
$\kappa\in\{101,211\}$ the full exact criterion at
$\varepsilon=\varepsilon_0(\kappa)$ holds with margins $180.6$ and $396.1$
(essentially the full main term).
\end{remark}

% ============================================================
\section{Off-Line Robustness}
\label{sec:offline}
% ============================================================

Sections~\ref{sec:id}--\ref{sec:uniform} are proved on the
prime/pole/archimedean side and assume nothing about zero locations.
We now isolate and bound the contribution of any hypothetical off-critical
zeros to the zero-side sum $\ZtwoGW=\Ztwoline+\Eoff$, where
$\Eoff(\varepsilon)=\tfrac12\sum_{\rho:\beta\ne 1/2}
h^{(2)}_{p,\varepsilon}\bigl((\rho-\tfrac12)/i\bigr)$.

\begin{lemma}[Quadruple form and $\beta$-uniform magnitude bound; proved]
\label{lem:mag}
Off-critical zeros occur in quadruples
$Q=\{\rho,\bar\rho,1-\rho,1-\bar\rho\}$ with $\rho=\beta+i\gamma$,
$\delta:=\beta-\tfrac12\ne 0$.
The quadruple contribution is real,
$\operatorname{contrib}_Q(\varepsilon)=2\,\Re\,h^{(2)}_{p,\varepsilon}(z_1)$
with $z_1=\gamma-i\delta$, and satisfies, independently of $\beta$,
\[
\bigl|\operatorname{contrib}_Q(\varepsilon)\bigr|
 \le 2\bigl(\gamma^2+\tfrac14\bigr)e^{\varepsilon^2/4}\sqrt p\;
   e^{-\varepsilon^2\gamma^2}.
\]
\end{lemma}

\begin{proof}
Evenness gives $h^{(2)}(-\gamma\mp i\delta)=h^{(2)}(\gamma\pm i\delta)$, and
real Taylor coefficients give
$h^{(2)}(\gamma+i\delta)=\overline{h^{(2)}(z_1)}$, so the quadruple sum
collapses to $2\Re h^{(2)}(z_1)$ after the zero-side factor $\tfrac12$.
Then $|z_1|^2=\gamma^2+\delta^2\le\gamma^2+\tfrac14$,
$|e^{-\varepsilon^2 z_1^2}|=e^{-\varepsilon^2(\gamma^2-\delta^2)}
\le e^{-\varepsilon^2\gamma^2}e^{\varepsilon^2/4}$, and
$|\cos(z_1\log p)|\le e^{|\delta|\log p}=p^{|\delta|}<\sqrt p$, using
$|\delta|<\tfrac12$.
\end{proof}

\begin{theorem}[Off-line robustness, magnitude form; proved]
\label{thm:offline}
Let $\kappa$ be the prime cutoff, $H_0=3\cdot10^{12}$, and
$f_\varepsilon(t):=(t^2+\tfrac14)e^{-\varepsilon^2 t^2}$, which is decreasing
on $[H_0,\infty)$ once $\varepsilon H_0\ge 2$. Then
\[
\sum_{p\le\kappa}(\log p)^2\bigl|\Eoff(\varepsilon)\bigr|
 \;\le\; 2\,C_\kappa\,e^{\varepsilon^2/4}\,
 S_{\mathrm{tail}}^{\mathrm{disc}}(\varepsilon,H_0),
\qquad
C_\kappa=\sum_{p\le\kappa}(\log p)^2\sqrt p,
\]
where the discrete tail
\[
S_{\mathrm{tail}}^{\mathrm{disc}}(\varepsilon,H_0)
 :=\sum_{n\ge 0}\bigl(N(H_0+n+1)-N(H_0+n)\bigr)\,f_\varepsilon(H_0+n)
\]
uses, for each bin count, the explicit Hasanalizade--Shen--Wong inequality
$|N(T)-F(T)|\le E(T)$ of~\cite{HasanalizadeShenWong21}, with
\[
F(T)=\tfrac{T}{2\pi}\log\tfrac{T}{2\pi e},
\qquad
E(T)=0.1038\log T+0.2573\log\log T+9.3675,
\]
so that each bin count satisfies the explicit estimate
\[
N(H_0+n+1)-N(H_0+n)\;\le\;F(H_0+n+1)-F(H_0+n)+E(H_0+n+1)+E(H_0+n).
\]
We use the Hasanalizade--Shen--Wong constants because they are sufficient for
the present threshold and match the archived certificate. Newer explicit
zero-count estimates of Bellotti--Wong~\cite{BellottiWong} improve these
constants and could slightly sharpen the numerical value of
$\varepsilon_{\mathrm{off}}$, but are not needed for the proof.
Consequently, for each certified negativity margin $M$ there is an explicit
threshold $\varepsilon_{\mathrm{off}}(\kappa,M,H_0)$ with
$\sum_{p\le\kappa}(\log p)^2|\Eoff(\varepsilon)|<M$ for all
$\varepsilon\ge\varepsilon_{\mathrm{off}}$.
Here $M$ is the \emph{scaled} interval-criterion margin of
Theorem~\ref{thm:ref} (e.g.\ $0.83$ or $21.86$); since
$8\sqrt\pi\varepsilon^3<1$ on $(0,0.05]$, it lies below the unscaled
negativity margin $-\BinfGW=(8\sqrt\pi\varepsilon^3)^{-1}\cdot(\text{scaled
gap})$, so $\sum_{p\le\kappa}(\log p)^2|\Eoff|<M$ is a conservative
sufficient condition for $\Binfline<0$.
For $\kappa=53$, $C_{53}=782.89$. The proved closed-form majorant above
\emph{exists} as an analytic threshold; its concrete decimal value is obtained
by numerical evaluation (mpmath, not Arb-interval-certified). Against
the smallest reference margin $M=0.83$ it gives
$\varepsilon_{\mathrm{off}}(53,0.83,H_0)\le 3.20\cdot10^{-12}$, and against the uniform
margin $M=S_3/4=21.86$,
$\varepsilon_{\mathrm{off}}(53,21.86,H_0)\le 3.14\cdot10^{-12}$
(these values are $\kappa=53$-specific, not uniform in $\kappa$, and
numerical); the independent numerical discrete evaluation gives
$3.15\cdot10^{-12}$ and $3.09\cdot10^{-12}$
respectively, the closed form agreeing to within $1\%$ (the gap is the
slowly-varying bin-density bound).
Hence, on a certified full-object negativity range with margin $M$ and for
$\varepsilon\ge\varepsilon_{\mathrm{off}}(\kappa,M,H_0)$ --- that is, for
$\varepsilon_{\mathrm{off}}(\kappa,M,H_0)\le\varepsilon\le 0.05$, in particular at
$\varepsilon=0.05$ (ten orders of magnitude above
$\varepsilon_{\mathrm{off}}$) --- the certified sign of $\BinfGW$ transfers
to the on-line proxy: $\Binfline(53,\varepsilon)<0$.
\end{theorem}

\begin{proof}
\emph{Step 1 (finite verified range).}
By Platt--Trudgian~\cite{PlattTrudgian21} there is no off-critical zero with
$0<\gamma\le H_0$ (the verified height is $3\,000\,175\,332\,800$, above
$3\cdot10^{12}$), so $\sum_p(\log p)^2|\Eoff|$ reduces to off-critical
quadruples with $\gamma>H_0$.

\emph{Step 2 (off-line support and discrete bin majorant).}
By the object split, $\Eoff=\ZtwoGW-\Ztwoline$ receives no contribution from
on-line zeros (there the full and proxy summands coincide), so $\Eoff$ is
supported on the off-critical quadruples $Q$ of Lemma~\ref{lem:mag}, each with
$|\operatorname{contrib}_Q|\le f_\varepsilon(\gamma)\sqrt p\,e^{\varepsilon^2/4}$.
Writing $N^{\mathrm{off}}$ for the off-critical counting function, and using
that the off-critical zeros in any bin are a subset of all zeros there
($N^{\mathrm{off}}(H_0+n+1)-N^{\mathrm{off}}(H_0+n)\le N(H_0+n+1)-N(H_0+n)$,
with $f_\varepsilon\ge0$),
\[
\sum_{\substack{\gamma>H_0\\ \beta\ne\tfrac12}} f_\varepsilon(\gamma)
 \le \sum_{n\ge 0}\bigl(N(H_0+n+1)-N(H_0+n)\bigr)\,f_\varepsilon(H_0+n)
 = S_{\mathrm{tail}}^{\mathrm{disc}}(\varepsilon,H_0),
\]
$f_\varepsilon$ being decreasing on $[H_0,\infty)$ for $\varepsilon H_0\ge2$.
This majorizes the off-line difference by the total zero count; the
over-count (on-line zeros contribute $0$ to $\Eoff$ but are counted in $N$) is
harmless for a magnitude statement and avoids any zero-density input.

\emph{Step 3 (zero-count bound).}
Each bin count $N(H_0+n+1)-N(H_0+n)$ is bounded by differencing the
explicit inequality of~\cite{HasanalizadeShenWong21}; this replaces the
Riemann--von Mangoldt density and is a rigorous upper bound for a discrete
zero sum. Multiplying by $2e^{\varepsilon^2/4}\sqrt p\,(\log p)^2$ and summing
over $p\le\kappa$ gives the stated bound.

\emph{Step 4 (log-density summation and threshold comparison).}
The map $\varepsilon\mapsto 2C_\kappa e^{\varepsilon^2/4}
S_{\mathrm{tail}}^{\mathrm{disc}}$ is strictly decreasing, and the threshold
$\varepsilon_{\mathrm{off}}$ admits a \emph{closed-form} majorant rather than
a bare infinite-series bisection. The bin count is \emph{not} uniformly
bounded: $D_n:=N(H_0+n+1)-N(H_0+n)=O(\log(H_0+n))$ grows with $n$. By the
explicit HSW inequality and the concavity bound
$\log(H_0+n)\le\log H_0+n/H_0$ one has the affine majorant
$D_n\le d_0+d_1 n$ with the \emph{secant} (not tangent) initial value
$d_0=F(H_0+1)-F(H_0)+E(H_0+1)+E(H_0)$,
$F(t)=\tfrac{t}{2\pi}\log\tfrac{t}{2\pi}-\tfrac{t}{2\pi}+\tfrac78$, and
$d_1=\tfrac1{2\pi H_0}+2E'(H_0)$, where
$E'(H_0)=\tfrac{0.1038}{H_0}+\tfrac{0.2573}{H_0\log H_0}$ captures the
logarithmic growth of the HSW error term over the bins, so
$d_1\approx\tfrac1{2\pi H_0}+\tfrac{0.21}{H_0}$ --- still negligible against
$H_0=3\cdot10^{12}$.
Keeping the full $f_\varepsilon(t)=(t^2+\tfrac14)e^{-\varepsilon^2t^2}$
(the $+\tfrac14$ must not be dropped from an upper bound) and using, for $n\ge0$,
\[
f_\varepsilon(H_0+n)\le \bigl((H_0+n)^2+\tfrac14\bigr)\, e^{-\varepsilon^2 H_0^2}\,q^{n},
\qquad q:=e^{-2\varepsilon^2 H_0}\in(0,1),
\]
(from $e^{-\varepsilon^2(H_0+n)^2}\le e^{-\varepsilon^2H_0^2}e^{-2\varepsilon^2H_0 n}$
and $e^{-\varepsilon^2 n^2}\le1$), the tail collapses to the geometric sums
$\sum q^n,\sum n q^n,\sum n^2 q^n,\sum n^3 q^n$ (all closed forms in $q$):
\[
S_{\mathrm{tail}}^{\mathrm{disc}}(\varepsilon,H_0)
 \le e^{-\varepsilon^2 H_0^2}
 \Bigl[d_0\bigl(\Sigma_2+\tfrac14\Sigma_0\bigr)
      +d_1\bigl(\Sigma_2'+\tfrac14\Sigma_1\bigr)\Bigr],
\]
with $\Sigma_0=\sum q^n$, $\Sigma_1=\sum nq^n$, $\Sigma_2=\sum(H_0+n)^2q^n$,
$\Sigma_2'=\sum n(H_0+n)^2q^n$, each a closed expression in $q$ evaluable in
ball arithmetic. The $\log$-growth term $d_1 n$ changes only the prefactor;
solving $2C_\kappa e^{\varepsilon^2/4}S_{\mathrm{tail}}^{\mathrm{disc}}=M$
yields the threshold directly.
At $\varepsilon=0.05$ the leading bin already carries the factor
$f_{0.05}(H_0)=O\bigl(H_0^2 e^{-0.0025\,H_0^2}\bigr)$, far beneath every
margin; the certified sign therefore transfers to $\Binfline$, and with
$O''(\tfrac12)=-2\Binfline$ the operator curvature is positive in that range.
\end{proof}

\begin{corollary}[Operator curvature positivity; proved]
\label{cor:curvline}
\emph{Reference case $\kappa=53$.} For every $\varepsilon$ with
$\varepsilon_{\mathrm{off}}(53,M,H_0)\le\varepsilon\le 0.05$ (in particular
$\varepsilon=0.05$, with the concrete thresholds
$\varepsilon_{\mathrm{off}}(53,0.83,H_0)\le3.20\cdot10^{-12}$ and
$\varepsilon_{\mathrm{off}}(53,21.86,H_0)\le3.14\cdot10^{-12}$, which are
numerical and not Arb-certified, and are not needed at sharp precision since
$\varepsilon=0.05$ lies far above them), the on-line
proxy curvature sum is negative, $\Binfline(53,\varepsilon)<0$, and hence the
operator curvature is strictly positive:
\[
O''(\tfrac12)=-2\Binfline(53,\varepsilon)>0.
\]
\emph{Uniform case $\kappa\ge 202$.} For every
$\varepsilon_{\mathrm{off}}(\kappa,M,H_0)\le\varepsilon\le\varepsilon_0(\kappa)$
(where nonempty) one has $\Binfline(\kappa,\varepsilon)<0$ and
$O''(\tfrac12)=-2\Binfline(\kappa,\varepsilon)>0$, by
Theorem~\ref{thm:uniform} and the transfer of Theorem~\ref{thm:offline}.
\end{corollary}

\begin{proof}
Theorems~\ref{thm:ref} and~\ref{thm:uniform} give $\BinfGW<0$ in the stated
ranges. For $\varepsilon\ge\varepsilon_{\mathrm{off}}$, Theorem~\ref{thm:offline}
bounds the off-line sum below the certified (scaled) margin $M$, which
satisfies $M<-\BinfGW$ on the certified range. Since
\[
\Binfline=\BinfGW-\sum_{p\le\kappa}(\log p)^2\Eoff,
\qquad
\Bigl|\sum_{p\le\kappa}(\log p)^2\Eoff\Bigr|<M<-\BinfGW=|\BinfGW|,
\]
the sign of $\Binfline$ equals that of $\BinfGW$, i.e.\ negative;
$O''(\tfrac12)=-2\Binfline$ then gives the curvature sign.
\end{proof}

\begin{remark}[What is and is not proved here]
Theorem~\ref{thm:offline} is a \emph{magnitude} statement: it shows that
hypothetical off-critical zeros leave the negativity, and hence
$O''(\tfrac12)>0$, intact for every $\varepsilon$ in a certified range with
$\varepsilon\ge\varepsilon_{\mathrm{off}}$ (that is,
$\varepsilon_{\mathrm{off}}\le\varepsilon\le 0.05$), in particular at the
reference width $\varepsilon=0.05$.
It does not control the \emph{sign} of $\Eoff$, and it does not prove RH.
Since the negativity theorems are already established on the zero-location-
free prime side, the distinct content is the explicit, $\beta$-uniform,
finite-height certificate quantifying the gap between the on-line picture
and the true zero configuration.
Subsequent work has refined the complementary zero-free and zero-density
regime above $H_0$ rather than the verified height itself
(\cite{MossinghoffTrudgianYang24,Bellotti24a}).
\end{remark}

% ============================================================
\section{Open Problems}
\label{sec:open}
% ============================================================

\begin{openproblem}[Cancellation Hypothesis]
\label{op:cancel}
Establish Assumption~(A$^{**}$) of~\cite{Paper7} analytically: a logarithmic
saving in $|S_k(\kappa)|/P(\kappa)$ via Vinogradov--Korobov exponential-sum
techniques (see~\cite{MossinghoffTrudgianYang24,Bellotti24a} for the state
of the art on zero-free regions).
Assumption~(A$^{**}$) is a log-saving hypothesis, strictly weaker than the
Riemann Hypothesis; no equivalence is claimed.
\end{openproblem}

\begin{openproblem}[Weighted Bias Bridge]
\label{op:bridge}
Establish the bridge implication
$B_{\mathrm{int}}^\infty<0\Rightarrow\BinfGW<0$ directly. The present paper
does not use it: it proves $\BinfGW<0$ independently and so bypasses the
bridge. Whether the integrated bias $B_{\mathrm{int}}^\infty$ controls the
full Guinand--Weil object through such an implication remains open.
\end{openproblem}

\begin{openproblem}[Exact stationarity]
\label{op:stat}
Establish $O'(\tfrac12)=0$, or $|O'(\tfrac12)|/O''(\tfrac12)\to 0$.
The selection criterion of~\cite{Paper7} quantifies the proximity to
stationarity but not its exactness; establishing exactness remains open.
\end{openproblem}

\begin{openproblem}[Weil-transfer operator and positivity]
\label{op:weil}
For the prime-local antisymmetric Gaussian component
$\psi_{p,\varepsilon}(x)=\phi_\varepsilon(x-\log p)-\phi_\varepsilon(x+\log p)$,
with $\phi_\varepsilon$ the Gaussian
$(2\pi\varepsilon^2)^{-1/2}\exp(-x^2/2\varepsilon^2)$, define the
Weil-transfer matrix
\[
(\mathsf{W}^{\mathrm{Weil}}_{\kappa,\varepsilon})_{pq}
:=W(\psi_{p,\varepsilon}\ast\check\psi_{q,\varepsilon}),
\]
where $W$ is the
Weil quadratic functional of~\cite{Lagarias,Bombieri}.
Establish whether a bridge identity relates
$\langle\mathsf{W}^{\mathrm{Weil}}_{\kappa,\varepsilon}c^\eta,c^\eta\rangle$,
with the coefficient vector $c^\eta=(\sqrt{V_p(\tfrac12)})_{p\le\kappa}$,
to $O''(\tfrac12)$ up to off-diagonal and $\Gamma$-factor corrections, and
whether the resulting form is positive on the growing finite-cutoff
subspace.
The vector $c^\eta$ is not asserted to reproduce the $(\log p)^2$ curvature
weights directly; identifying the exact renormalisation and correction terms
between the $c^\eta$-weighted Weil form and the curvature weight is part of
the problem.
Positivity of the Weil form on this subspace would be a finite-cutoff
analogue of the Weil positivity criterion of~\cite{Lagarias}, which is known
to be equivalent to the Riemann Hypothesis~\cite{Lagarias,Bombieri}.
\end{openproblem}

\begin{openproblem}[Scaling of the negativity window]
\label{op:scaling}
Determine the asymptotics of the negativity threshold; in particular,
whether $\inf_\kappa\varepsilon_{\mathrm{int}}(\kappa)>0$ for the integrated
window of~\cite{Paper7}, and how the uniform threshold $\varepsilon_0(\kappa)$
of Theorem~\ref{thm:uniform} compares with the observed scale
$\dmin(\kappa)\asymp 2/\kappa$.
Numerical evidence is oscillatory.
\end{openproblem}

\begin{remark}[Scope: the curvature is not a direct RH criterion]
The negativity $\BinfGW<0$ established here is a statement about the
second-moment bias, and the operator curvature
$O''(\tfrac12)=-2\Binfline>0$ follows by the off-line transfer; neither is a
criterion for the Riemann Hypothesis.
The canonical off-critical continuation of the weighted curvature test
function
\[
h^{(2,w)}_{p,\varepsilon}(u):=(\log p)^2 h^{(2)}_{p,\varepsilon}(u)
=(\log p)^2 u^2 e^{-\varepsilon^2 u^2}\cos(u\log p),
\]
written $G(\gamma):=h^{(2,w)}_{p,\varepsilon}(\gamma)$, is obtained as follows.
A hypothetical off-critical zero $\rho=\tfrac12+\delta+i\gamma$ enters the
half-sum $B=\tfrac12\sum_\rho$ together with its complex conjugate
$\overline\rho=\tfrac12+\delta-i\gamma$, contributing the arguments
$u=(\rho-\tfrac12)/i=\gamma-i\delta$ and $u=(\overline\rho-\tfrac12)/i=-\gamma-i\delta$;
since $h^{(2,w)}_{p,\varepsilon}$ is even in $u$, $G(-\gamma-i\delta)=G(\gamma+i\delta)$,
so this conjugate pair contributes $G(\gamma+i\delta)+G(\gamma-i\delta)$ to
$\sum_\rho$, i.e.\ $\tfrac12\bigl[G(\gamma+i\delta)+G(\gamma-i\delta)\bigr]$ in the
half-sum, with on-line limit $G(\gamma)$ at $\delta=0$. The defect relative to
that limit is
\[
\Delta B(\delta,\gamma)
:=\tfrac12\bigl[G(\gamma+i\delta)+G(\gamma-i\delta)-2G(\gamma)\bigr]
=\tfrac12\bigl[-\delta^2 G''(\gamma)+O(\delta^4)\bigr]
=\delta^2 C_2(\gamma)+O(\delta^4),
\]
the middle step being the even-order Taylor expansion
$G(\gamma\pm i\delta)=G(\gamma)\mp\tfrac{\delta^2}{2}G''(\gamma)+O(\delta^4)$ along
the imaginary direction. Hence
\[
C_2(\gamma)=-\tfrac12 G''(\gamma),
\]
the factor $-\tfrac12$ originating in the half-sum $\tfrac12\sum_\rho$; the
functional-equation partners $1-\rho,1-\overline\rho$ form an identical conjugate
pair contributing the same defect, and the un-normalised convention
$\Delta B:=G(\gamma+i\delta)+G(\gamma-i\delta)-2G(\gamma)$ (without the half-sum
$\tfrac12$) would instead read off $-G''(\gamma)$. The coefficient $C_2(\gamma)$ is
sign-indefinite: for $p=2$, $\varepsilon=0.05$, $G''(\gamma)$ changes sign $13$
times on $[1,60]$ (e.g.\ $G''(5)\approx+6.41>0$ but $G''(9)\approx-14.85<0$),
consistent with the larger-range numerical scans; equivalently, the
Fourier transform of a second derivative is not of fixed sign.
Hence the curvature bias does not furnish a Weil-positivity / RH criterion.
This delimits the result and is stated for transparency.
\end{remark}

% ============================================================
\section*{Non-Circularity}
\label{sec:nocirc}
% ============================================================

We document explicitly that the results of this paper do not rely on the
objects they are directed toward.

\textbf{No RH as input.}
The Riemann Hypothesis is never assumed.
All zero sums are defined through the zero side of the explicit formula and
require no hypothesis on zero positions; the ordinate form quoted in
\S\,\ref{sec:bg} coincides with this definition on the range of zeros
entering any numerical computation.
The results of \S\S\,\ref{sec:id}--\ref{sec:uniform} are unconditional, and
\S\,\ref{sec:offline} is a magnitude statement.

\textbf{No GUE, Montgomery, or Hilbert--P\'olya hypothesis.}
The distribution of the zeros is not modelled by random matrix theory.
The operator $\Tt$ is not a Hilbert--P\'olya operator (\cite{Paper4}:
$r_2\to 0.16$); its matrix entries are explicit deterministic arithmetic
expressions.
No pair-correlation conjecture is used.

\textbf{Off-line robustness uses a verified height, not RH.}
The single external input of \S\,\ref{sec:offline} is the Platt--Trudgian
computation~\cite{PlattTrudgian21}, a rigorous unconditional
interval-arithmetic verification that no off-critical zero exists below
$H_0=3\cdot10^{12}$.
This is a finite arithmetic fact, not an instance of RH.

\textbf{$\BinfGW<0$ is not assumed and not RH-equivalent.}
The negativity is established by closed-form and certified computation
(Theorems~\ref{thm:closedform},~\ref{thm:ref},~\ref{thm:uniform}); it is a
one-sided curvature/bias fact, and by the scope remark of \S\,\ref{sec:open}
it is not a criterion for RH.
The non-vanishing $\zeta(1+it)\ne 0$~\cite{Hadamard,dlVP} is an
unconditional analytical input.

% ============================================================
\section*{Acknowledgments}
% ============================================================

The author thanks the readers who provided feedback on earlier versions of
this paper and on the preceding papers in the series.

% ============================================================
\section*{Call for Feedback}
% ============================================================

This paper is part of an open research initiative on prime--zero
coupling and explicit-formula operators.
All papers in the series are freely available on Zenodo under CC~BY~4.0, and
the accompanying verification code is hosted on GitHub.

{\raggedright
Zenodo:\\
\url{https://zenodo.org/search?q=metadata.creators.person_or_org.name\%3A\%22Tehrani\%2C\%20Ulrich\%22}\\
GitHub:\\
\url{https://github.com/utehrani/analysislab-nt}
\par}

Topics of particular interest include:
analytical approaches to Assumption~(A$^{**}$) via exponential-sum methods;
the relation between the negativity window and zero-free regions for
$\zeta(s)$; sharper archimedean constants; and any errors or gaps in the
present arguments.

% ============================================================
\newpage
\begin{thebibliography}{99}

\bibitem{Paper6}
U.~Tehrani,
``Positive Curvature of the Spectral Trace at the Critical Line'',
Zenodo, \url{https://doi.org/10.5281/zenodo.19665790}, 2026.

\bibitem{Paper5}
U.~Tehrani,
``Spectral Trace Formula and Smoothed Zero Sums:
A Prime--Zero Duality Framework'',
Zenodo, \url{https://doi.org/10.5281/zenodo.19508547}, 2026.

\bibitem{Paper7}
U.~Tehrani,
``Conditional Stationarity and Positive Curvature of the Spectral Trace
at the Critical Line'',
Zenodo, \url{https://doi.org/10.5281/zenodo.20440671}, 2026.

\bibitem{Code}
U.~Tehrani,
Verification scripts for the prime--zero coupling series,
GitHub, \url{https://github.com/utehrani/analysislab-nt}.

\bibitem{Paper1}
U.~Tehrani,
``A Curvature Decomposition of the Explicit Formula
for the Riemann Zeta Function'',
Zenodo, \url{https://doi.org/10.5281/zenodo.19025598}, 2026.

\bibitem{Paper2}
U.~Tehrani,
``From Local Curvature to the Weil Functional:
An Explicit Construction'',
Zenodo, \url{https://doi.org/10.5281/zenodo.19106992}, 2026.

\bibitem{Lagarias}
J.C.~Lagarias,
``On a positivity property of the Riemann $\xi$-function'',
\emph{Acta Arithmetica} \textbf{89} (1999), 217--234.

\bibitem{Paper3}
U.~Tehrani,
``A Finite-Cutoff Hilbert-Space Model for Prime--Zero Energy Structure'',
Zenodo, \url{https://doi.org/10.5281/zenodo.19307989}, 2026.

\bibitem{Paper4}
U.~Tehrani,
``A Dual Operator for Prime--Zero Coupling
and a Conditional Proof of Energy Asymmetry'',
Zenodo, \url{https://doi.org/10.5281/zenodo.19364703}, 2026.

\bibitem{BerryKeating}
M.V.~Berry and J.P.~Keating,
``The Riemann Zeros and Eigenvalue Asymptotics'',
\emph{SIAM Review} \textbf{41}.2 (1999), 236--266.

\bibitem{deBranges}
L.~de~Branges,
``The Riemann Hypothesis for Hilbert Spaces of Entire Functions'',
\emph{Bull.\ Amer.\ Math.\ Soc.\ (N.S.)} \textbf{15}.1 (1986), 1--17.

\bibitem{ConnesConsani23}
A.~Connes and C.~Consani,
``Spectral triples and zeta-cycles,''
\textit{Enseign.\ Math.} \textbf{69} (2023), no.~1--2, 93--148.
arXiv:2106.01715.

\bibitem{CCM25}
A.~Connes, C.~Consani, and H.~Moscovici,
``Zeta spectral triples,''
arXiv:2511.22755 (2025).

\bibitem{Bombieri}
E.~Bombieri,
``Remarks on Weil's Quadratic Functional in the Theory of Prime
Numbers, I'',
\emph{Rend.\ Lincei Mat.\ Appl.} \textbf{11}.3 (2000), 183--233.

\bibitem{Burnol}
J.-F.~Burnol,
``On Fourier and Zeta(s)'',
\emph{Forum Math.} \textbf{16}.6 (2004), 789--840.

\bibitem{Suzuki25}
M.~Suzuki,
``On the Hilbert space derived from the Weil distribution,''
\textit{Canad.\ J.\ Math.}, First View (2025), 1--24.
DOI: 10.4153/S0008414X25101739.
arXiv:2301.00421.

\bibitem{Guinand}
A.P.~Guinand,
``A summation formula in the theory of prime numbers'',
\emph{Proc.\ London Math.\ Soc.}\ (2) \textbf{50} (1948), 107--119.

\bibitem{Weil}
A.~Weil,
``Sur les `formules explicites' de la th\'{e}orie des nombres premiers'',
\emph{Comm.\ S\'{e}m.\ Math.\ Univ.\ Lund, Suppl.}\ (1952), 252--265.

\bibitem{IwaniecKowalski04}
H.~Iwaniec and E.~Kowalski,
\emph{Analytic Number Theory},
American Mathematical Society Colloquium Publications,
vol.~53, AMS, Providence, RI, 2004.

\bibitem{BalanzarioCardenas23}
E.P.~Balanzario and D.E.~C\'{a}rdenas Romero,
``An explicit formula for the zeros of the Riemann zeta function
and the statistics of their local distribution'',
\emph{The Ramanujan Journal} \textbf{69} (2026), Article~21,
DOI: 10.1007/s11139-025-01297-y; arXiv:2312.00108.

\bibitem{RosserSchoenfeld62}
J.B.~Rosser and L.~Schoenfeld,
``Approximate formulas for some functions of prime numbers'',
\emph{Illinois J.\ Math.} \textbf{6} (1962), 64--94.

\bibitem{HasanalizadeShenWong21}
E.~Hasanalizade, Q.~Shen, and P.-J.~Wong,
``Counting zeros of the Riemann zeta function'',
\emph{J.\ Number Theory} \textbf{235} (2022), 219--241.
arXiv:2107.06506.

\bibitem{BellottiWong}
C.~Bellotti and P.-J.~Wong,
``Improved estimates for the argument and zero-counting function of the
Riemann zeta-function'' (with an appendix by A.~Fiori),
\emph{Math.\ Comp.} (2025), DOI: 10.1090/mcom/4133; arXiv:2412.15470.

\bibitem{PlattTrudgian21}
D.J.~Platt and T.~Trudgian,
``The Riemann hypothesis is true up to $3\cdot10^{12}$'',
\emph{Bull.\ Lond.\ Math.\ Soc.} \textbf{53} (2021), no.~3, 792--797.

\bibitem{BrentPlattTrudgian20}
R.P.~Brent, D.J.~Platt, and T.S.~Trudgian,
``The mean square of the error term in the prime number theorem'',
\emph{J.\ Number Theory} \textbf{238} (2022), 740--762.
arXiv:2010.06451.

\bibitem{MossinghoffTrudgianYang24}
M.J.~Mossinghoff, T.S.~Trudgian, and A.~Yang,
``Explicit zero-free regions for the Riemann zeta-function,''
\textit{Res.\ Number Theory} \textbf{10} (2024), art.~11.
DOI: 10.1007/s40993-023-00498-y.

\bibitem{Bellotti24a}
C.~Bellotti,
``Explicit bounds for the Riemann zeta function
and a new zero-free region,''
\textit{J.\ Math.\ Anal.\ Appl.} \textbf{536} (2024),
no.~2, art.~128249.
DOI: 10.1016/j.jmaa.2024.128249.
arXiv:2306.10680.

\bibitem{Johansson17}
F.~Johansson,
``Arb: efficient arbitrary-precision midpoint-radius interval arithmetic,''
\emph{IEEE Trans.\ Comput.} \textbf{66}.8 (2017), 1281--1292.
arXiv:1611.02831.

\bibitem{Hadamard}
J.~Hadamard,
``Sur la distribution des z\'{e}ros de la fonction $\zeta(s)$
et ses cons\'{e}quences arithm\'{e}tiques'',
\emph{Bull.\ Soc.\ Math.\ France} \textbf{24} (1896), 199--220.

\bibitem{dlVP}
C.-J.~de~la~Vall\'{e}e~Poussin,
``Recherches analytiques sur la th\'{e}orie des nombres premiers'',
\emph{Ann.\ Soc.\ Sci.\ Bruxelles} \textbf{20} (1896), 183--256.

\bibitem{AbramowitzStegun}
M.~Abramowitz and I.A.~Stegun,
\emph{Handbook of Mathematical Functions},
National Bureau of Standards, Applied Mathematics Series \textbf{55},
Washington, DC, 1964; \S6.3 (psi function).

\end{thebibliography}
% ============================================================

\enlargethispage{3\baselineskip}
\vspace*{\fill}
\begin{minipage}{\linewidth}
\rule{\linewidth}{0.4pt}\smallskip\\
\small\emph{Paper~8 v0.23 \textperiodcentered\ August~2026
\textperiodcentered\ Pauca sed matura}
\end{minipage}

\end{document}
